% !TEX root = TA.tex
% ==========================================
% BAB VI EVALUASI
% ==========================================
\cleardoublepage
\chapter{EVALUASI}
\label{chap:evaluasi}

Bab ini membahas hasil pengujian pengaruh persona terhadap respons enam
model LLM. Rancangan eksperimen dan proses pengujian masing-masing telah
dijelaskan pada Bab~\ref{chap:perancangan} dan
Bab~\ref{chap:implementasi}. Evaluasi dilakukan pada GSM8K dan ETHICS untuk
domain gender, agama, dan regional-kultural. Pembahasan mencakup akurasi dan
penanda linguistik.

% =====================================================================
\section{Evaluasi GSM8K}
\label{sec:evaluasi-gsm8k}

GSM8K terdiri atas 1.319 soal matematika. Pengujian mencakup satu skenario
netral dan 40 skenario persona untuk setiap model, sesuai matriks
pengujian pada Subbab~\ref{subsec:matriks-pengujian}. Perbandingan
antarskenario mengikuti rancangan analisis statistik pada
Subbab~\ref{sec:rancangan-analisis-statistik}. Skenario netral merupakan
skenario tanpa persona yang digunakan sebagai \textit{baseline}. Akurasi
dihitung sebagai persentase respons yang menghasilkan angka GSM8K sesuai
dengan \textit{ground truth} pada dataset.

Hasil pada skenario netral menjadi pembanding dan ditampilkan sebagai
baris N\label{acr:n} pada setiap tabel domain. Akurasi pada skenario ini
berkisar antara 84,43\% dan 94,08\%. Qwen3 Next 80B-A3B memperoleh
akurasi tertinggi, sedangkan Llama 3.1 8B memperoleh akurasi terendah.
Empat model lainnya memperoleh akurasi 91,96--93,93\%.

% =====================================================================
\subsection{Domain Gender}
\label{subsec:hasil-gender-gsm8k}

Domain gender mencakup persona laki-laki dan perempuan. Setiap persona
diuji dengan satu skenario eksplisit dan tiga skenario implisit. I1\label{acr:i1}
menggunakan peran dalam keluarga, I2\label{acr:i2} menggunakan kegiatan
merapikan penampilan, dan I3\label{acr:i3} menggunakan pilihan pakaian.
Ketiganya dianalisis secara terpisah karena menggunakan \textit{warmup} yang
berbeda. Analisis ini menggunakan 71.226 respons dari 1.319 soal GSM8K.

\subsubsection{Akurasi Respons}
\label{subsubsec:akurasi-gender-gsm8k}

Analisis akurasi domain gender membandingkan hasil pada
skenario netral, eksplisit, dan implisit. Akurasi E dan I1--I3 dibandingkan
dengan N, I1--I3 dibandingkan dengan E, serta persona laki-laki dan perempuan
dibandingkan pada soal dan skenario yang sama. Perbedaan antara N, E, I1, I2,
dan I3 diperiksa menggunakan Q Cochran, sedangkan setiap pasangan diuji
menggunakan Exact McNemar dengan koreksi Holm; hasil antarskenario dimasukkan
ke hasil utama jika Q Cochran pada kelompok yang sama juga signifikan.
Pengujian mencakup 108 perbandingan, yaitu 48 terhadap N, 36 terhadap E, dan 24
antara persona laki-laki dan perempuan. \textit{Input token},
\textit{completion token}, dan total token dihitung secara terpisah untuk
membandingkan jumlah token yang digunakan setiap persona dan skenario.

\par\noindent\textbf{a. Akurasi Setiap Skenario}\par

Tabel~\ref{tab:akurasi-gender-gsm8k} menunjukkan akurasi keenam model
pada skenario netral serta skenario yang menggunakan persona laki-laki
dan perempuan. Baris N berisi hasil pada skenario netral. Untuk setiap
persona, hasil skenario eksplisit ditulis pada baris E\label{acr:e},
sedangkan hasil tiga skenario implisit ditulis pada baris I1, I2, dan I3.

\begin{table}[H]
    \centering
    \caption{Akurasi GSM8K pada domain gender menurut persona dan skenario (\%)}
    \label{tab:akurasi-gender-gsm8k}
    \scriptsize
    \setlength{\tabcolsep}{3pt}
    \renewcommand{\arraystretch}{1.10}
    \begin{tabularx}{\textwidth}{
        >{\raggedright\arraybackslash}p{1.75cm}
        >{\centering\arraybackslash}p{0.65cm}
        *{6}{>{\centering\arraybackslash}X}}
        \toprule
        \textbf{Persona} & \textbf{Skenario} & \textbf{Qwen3} &
        \textbf{Llama 3.1} & \textbf{Ministral} & \textbf{DeepSeek} &
        \textbf{GPT-OSS} & \textbf{Llama 3.3} \\
        \midrule
        Netral & N & 94,08 & 84,43 & 93,47 & 93,93 & 93,93 & 91,96 \\
        \addlinespace[2pt]
        Laki-laki & E  & 92,65 & 83,24 & 93,40 & 93,78 & 94,77 & 92,49 \\
                  & I1 & 92,12 & 82,11 & 92,42 & 90,83 & 95,00 & 93,86 \\
                  & I2 & 92,57 & 84,84 & 93,18 & 91,66 & 95,15 & 93,48 \\
                  & I3 & 91,96 & 83,47 & 93,56 & 92,57 & 94,09 & 92,95 \\
        \addlinespace[2pt]
        Perempuan & E  & 93,71 & 82,64 & 93,18 & 93,40 & 94,69 & 91,89 \\
                  & I1 & 92,34 & 82,64 & 93,25 & 90,22 & 95,38 & 93,71 \\
                  & I2 & 92,71 & 84,76 & 93,33 & 91,05 & 94,84 & 93,71 \\
                  & I3 & 92,72 & 84,31 & 93,70 & 92,49 & 95,22 & 94,01 \\
        \bottomrule
    \end{tabularx}
\end{table}

Dalam satu model, akurasi antarskenario tidak jauh berbeda. Pada keenam
model, selisih antara akurasi tertinggi dan terendah dalam satu model
berkisar 1,28--3,71 poin persentase. Selisih paling besar terdapat pada
DeepSeek R1-Distill 32B, sedangkan selisih paling kecil terdapat pada
Ministral 8B. Gambar~\ref{fig:ringkasan-akurasi-gender-gsm8k}
menunjukkan perbandingan hasil pada skenario netral, skenario eksplisit,
dan rata-rata I1, I2, dan I3.

\begin{figure}[H]
    \centering
    \includegraphics[
        width=\textwidth,
        height=0.46\textheight,
        keepaspectratio
    ]{image/bab6/gsm8k/gender/akurasi/01_gender_absolute_accuracy_five_categories.png}
    \caption[Akurasi GSM8K menurut skenario pada domain gender]
    {Perbandingan akurasi GSM8K antara skenario netral, skenario eksplisit,
    dan rata-rata skenario implisit pada domain gender}
    \label{fig:ringkasan-akurasi-gender-gsm8k}
\end{figure}

Perbandingan dengan skenario netral membentuk tiga pola. Seluruh skenario
persona lebih rendah pada Qwen3 Next 80B-A3B dan DeepSeek R1-Distill 32B,
tetapi lebih tinggi pada GPT-OSS 20B. Pada tiga model lainnya, hasil yang lebih
tinggi hanya ditemukan pada sebagian skenario.

Gambar~\ref{fig:delta-gender-implisit-eksplisit-gsm8k} membandingkan rata-rata
I1--I3 dengan E untuk persona laki-laki dan perempuan serta membandingkan hasil
kedua persona. Pada batang biru dan merah, nilai positif menunjukkan rata-rata
I1--I3 lebih tinggi daripada E, sedangkan nilai negatif menunjukkan hasil yang
lebih rendah. Batang ungu dihitung dari E dan rata-rata I1--I3 setiap persona;
nilai positif menunjukkan persona laki-laki lebih tinggi dan nilai negatif
menunjukkan persona perempuan lebih tinggi.

\begin{figure}[H]
    \centering
    \includegraphics[
        width=\textwidth,
        height=0.46\textheight,
        keepaspectratio
    ]{image/bab6/gsm8k/gender/akurasi/03_gender_delta_implicit_vs_explicit.png}
    \caption[Selisih akurasi antarskenario dan antarpersona pada domain gender]
    {Selisih rata-rata akurasi antara skenario implisit dan eksplisit serta
    antara persona laki-laki dan perempuan pada domain gender}
    \label{fig:delta-gender-implisit-eksplisit-gsm8k}
\end{figure}

Pada DeepSeek R1-Distill 32B, rata-rata akurasi I1--I3 lebih rendah 2,08 poin
untuk persona laki-laki dan 2,15 poin untuk persona perempuan dibandingkan
dengan E. Sebagai pembanding, rata-rata akurasi I1--I3 pada Llama 3.3 70B
lebih tinggi 0,94 dan 1,92 poin. Rata-rata skenario implisit karena itu dapat
lebih rendah atau lebih tinggi daripada E menurut model yang digunakan.

\par\noindent\textbf{b. Hasil Perbandingan Statistik}\par
\label{sym:p-value}

Uji Q Cochran signifikan pada tiga dari 12 kelompok. Hasil tersebut ditemukan
pada kedua persona DeepSeek R1-Distill 32B dengan $p<0{,}001$ serta persona
perempuan Llama 3.3 70B dengan $p=0{,}012$. Setelah koreksi Holm, delapan dari
108 perbandingan tetap signifikan. Tujuh hasil melibatkan I1 atau I2 pada
DeepSeek, sedangkan satu hasil terdapat pada I3 persona perempuan Llama 3.3.

Empat hasil signifikan terhadap skenario netral seluruhnya berupa penurunan
pada I1 dan I2 DeepSeek R1-Distill 32B untuk kedua persona. Seluruh
perbandingan E dengan N tidak signifikan. Terhadap E, tiga penurunan juga
terdapat pada DeepSeek, sedangkan satu kenaikan terdapat pada I3 persona
perempuan Llama 3.3 70B. Besar perubahan setiap perbandingan dicantumkan pada
Tabel~\ref{tab:hasil-signifikan-gender-gsm8k}.

Dari 24 perbandingan antarpersona laki-laki dan perempuan, tidak ada
yang signifikan setelah koreksi Holm. Selisih paling besar hanya 1,14 poin
persentase dengan nilai $|h|$ paling besar 0,051. Hasil akurasi karena itu tidak
mendukung keunggulan salah satu persona gender.

Gambar~\ref{fig:effect-gender-gsm8k} menunjukkan \textit{effect size}
Cohen's $h$ dalam dua panel. Panel pertama membandingkan E dan I1--I3 dengan N,
sedangkan panel kedua membandingkan I1--I3 dengan E serta persona laki-laki
dengan perempuan. Warna biru menunjukkan akurasi skenario pertama lebih rendah
daripada pembanding, sedangkan warna merah menunjukkan akurasi yang lebih
tinggi. Tanda bintang menunjukkan hasil yang tetap signifikan setelah koreksi
Holm, sedangkan nilai \textit{p-value} lengkap terdapat pada
Lampiran~\ref{chap:lampiran-statistik-gender}.

\begin{figure}[H]
    \centering
    \includegraphics[
        width=\textwidth,
        height=0.78\textheight,
        keepaspectratio
    ]{image/bab6/gsm8k/gender/akurasi/24_gender_effect_holm_single_large_heatmap.png}
    \caption[Ukuran efek dan signifikansi akurasi pada domain gender]
    {Ukuran efek Cohen's $h$ dan hasil uji Exact McNemar setelah koreksi Holm
    pada 18 jenis perbandingan per model dalam domain gender.\\[-2pt]
    {\footnotesize\textit{Catatan:} * = signifikan setelah koreksi Holm
    ($p_{\mathrm{Holm}} < 0{,}05$); E = eksplisit; I1--I3 = implisit;
    N = netral; L = laki-laki; P = perempuan. Cohen's $h$ dihitung dari
    skenario pertama dikurangi skenario pembanding.}}
    \label{fig:effect-gender-gsm8k}
\end{figure}

Tabel~\ref{tab:hasil-signifikan-gender-gsm8k} memuat rincian delapan
perbandingan yang bertanda bintang pada gambar. Nilai $\Delta$, interval
kepercayaan, dan \textit{p-value} menunjukkan besar perbedaan dan hasil
pengujiannya. Cohen's $h$ menunjukkan \textit{effect size}, sedangkan MOR
menunjukkan perubahan jawaban pada soal yang sama.

\begin{table}[H]
    \centering
    \caption{Perbandingan domain gender yang tetap signifikan setelah koreksi Holm}
    \label{tab:hasil-signifikan-gender-gsm8k}
    \scriptsize
    \renewcommand{\arraystretch}{1.18}
    \setlength{\tabcolsep}{3pt}
    \begin{tabularx}{\textwidth}{
        >{\raggedright\arraybackslash}p{1.65cm}
        >{\raggedright\arraybackslash}X
        >{\centering\arraybackslash}p{0.85cm}
        >{\centering\arraybackslash}p{2.05cm}
        >{\centering\arraybackslash}p{1.05cm}
        >{\centering\arraybackslash}p{0.75cm}
        >{\centering\arraybackslash}p{0.75cm}}
        \toprule
        \textbf{Model} & \textbf{Perbandingan} &
        \textbf{$\Delta$} & \textbf{IK 95\%}\label{acr:ik} &
        \textbf{$p_{\mathrm{Holm}}$} & \textbf{$h$} &
        \textbf{MOR}\label{acr:mor} \\
        \midrule
        DeepSeek & Laki-laki I1--netral & $-3{,}03$ &
        $[-4{,}63;\,-1{,}44]$ & 0,003 & $-0{,}12$ & 0,51 \\
        DeepSeek & Laki-laki I2--netral & $-2{,}28$ &
        $[-3{,}87;\,-0{,}68]$ & 0,034 & $-0{,}09$ & 0,59 \\
        DeepSeek & Perempuan I1--netral & $-3{,}72$ &
        $[-5{,}39;\,-2{,}05]$ & \mbox{$<0{,}001$} & $-0{,}14$ & 0,44 \\
        DeepSeek & Perempuan I2--netral & $-2{,}89$ &
        $[-4{,}56;\,-1{,}21]$ & 0,004 & $-0{,}11$ & 0,53 \\
        DeepSeek & Laki-laki I1--eksplisit & $-2{,}96$ &
        $[-4{,}70;\,-1{,}29]$ & 0,005 & $-0{,}11$ & 0,55 \\
        DeepSeek & Perempuan I1--eksplisit & $-3{,}18$ &
        $[-4{,}78;\,-1{,}59]$ & \mbox{$<0{,}001$} & $-0{,}12$ & 0,48 \\
        DeepSeek & Perempuan I2--eksplisit & $-2{,}35$ &
        $[-4{,}02;\,-0{,}76]$ & 0,021 & $-0{,}09$ & 0,58 \\
        Llama 3.3 & Perempuan I3--eksplisit & $+2{,}12$ &
        $[+0{,}83;\,+3{,}49]$ & 0,018 & $+0{,}08$ & 2,04 \\
        \bottomrule
    \end{tabularx}

    \vspace{0.15cm}
    \begin{minipage}{0.98\textwidth}
        \scriptsize
        Nilai $\Delta$ dan interval kepercayaan ditulis dalam poin
        persentase. MOR adalah \textit{matched odds ratio}.
    \end{minipage}
\end{table}

Nilai $|h|$ pada delapan perbandingan signifikan berkisar 0,08--0,14
dan seluruhnya berada di bawah 0,20. Pada DeepSeek R1-Distill 32B, MOR sebesar
0,44--0,59 menunjukkan lebih banyak pasangan dengan jawaban benar pada
skenario pembanding dan jawaban salah pada skenario implisit. Pada Llama
3.3 70B, MOR sebesar 2,04 menunjukkan lebih banyak jawaban berubah dari
salah pada E menjadi benar pada I3.

\par\noindent\textbf{c. Penggunaan Token}\par

Skenario implisit menggunakan lebih banyak \textit{input token} karena seluruh
percakapan \textit{warmup} dikirim bersama soal. Rata-ratanya meningkat
213,7--965,8\% dibandingkan dengan E. I1 menggunakan \textit{input token}
paling banyak pada seluruh model dan kedua persona karena percakapan
\textit{warmup}-nya paling panjang. Tabel
\ref{tab:perubahan-token-gender} merangkum perubahan
\textit{completion token} pada setiap model.

\begin{table}[H]
    \centering
    \caption[Perubahan completion token pada domain gender]
    {Perubahan rata-rata \textit{completion token} I1--I3 terhadap E pada
    domain gender}
    \label{tab:perubahan-token-gender}
    \small
    \renewcommand{\arraystretch}{1.15}
    \begin{tabular}{lcc}
        \toprule
        \textbf{Model} & \textbf{Laki-laki} & \textbf{Perempuan} \\
        \midrule
        Qwen3 Next 80B-A3B      & $+10{,}0\%$ & $+10{,}3\%$ \\
        Llama 3.1 8B            & $+5{,}0\%$  & $-1{,}9\%$  \\
        Ministral 8B            & $-0{,}9\%$  & $-1{,}8\%$  \\
        DeepSeek R1-Distill 32B & $+82{,}8\%$ & $+80{,}0\%$ \\
        GPT-OSS 20B             & $-7{,}4\%$  & $-4{,}7\%$  \\
        Llama 3.3 70B           & $+4{,}3\%$  & $+3{,}3\%$  \\
        \bottomrule
    \end{tabular}

    \vspace{0.15cm}
    \begin{minipage}{0.88\textwidth}
        \scriptsize
        Persentase dihitung dari rata-rata I1--I3 dibandingkan dengan E.
    \end{minipage}
\end{table}

Perubahan \textit{completion token} berbeda antarmodel. DeepSeek
R1-Distill 32B mengalami kenaikan 80,0--82,8\%, sedangkan GPT-OSS 20B
mengalami penurunan 4,7--7,4\%. Pada empat model lainnya, perubahannya
berkisar $-1{,}9\%$ sampai $+10{,}3\%$. Seluruh model menerima tambahan
percakapan \textit{warmup}, tetapi panjang respons yang dihasilkan dapat
meningkat atau menurun.

Total token pada rata-rata I1--I3 meningkat 59,5--220,5\% terhadap E pada
seluruh model. Persona laki-laki menggunakan total token 0,6--26,7\% lebih
banyak daripada persona perempuan karena percakapan \textit{warmup}-nya lebih
panjang. \textit{Completion token} persona laki-laki tidak selalu lebih
tinggi, sehingga perbedaan jumlah \textit{input token} menjadi sumber utama
selisih total token.

Penggunaan total token yang lebih besar pada persona laki-laki tidak disertai
keunggulan akurasi yang signifikan. Seluruh perbandingan akurasi antarpersona
laki-laki dan perempuan memperoleh
$p_{\mathrm{Holm}}\geq 0{,}05$. Persona gender mengubah jumlah token yang
digunakan tanpa menghasilkan perbedaan akurasi antarpersona yang signifikan.
Rata-rata komponen token terdapat pada
Gambar~\ref{fig:lampiran-token-gender-components}, sedangkan perbandingan
token dan akurasi setiap skenario terdapat pada
Gambar~\ref{fig:lampiran-token-akurasi-gender} di
Lampiran~\ref{sec:lampiran-token-gender}.

\subsubsection{Penanda Linguistik}
\label{subsubsec:penanda-gender-gsm8k}

Perbandingan penanda linguistik menggunakan skenario dan pasangan persona yang
sama dengan pengujian akurasi. Pada GSM8K, perubahan penanda linguistik gender
terutama terdapat pada bentuk respons. Perubahannya mencakup panjang respons,
variasi kata, susunan penyelesaian, bahasa matematis, dan kata ganti diri.
Penanda linguistik yang berubah pada E, I1, I2, dan I3 berbeda menurut model. Perbandingan
antarpersona laki-laki dan perempuan juga menghasilkan perubahan
yang berbeda, sehingga pengaruh persona gender muncul melalui cara setiap model
menulis respons.

\par\noindent\textbf{a. Perubahan terhadap Skenario Netral}\par

Gambar~\ref{fig:marker-gender-terhadap-netral} membandingkan skenario E, I1,
I2, dan I3 dengan skenario netral pada persona yang sama. Baris atas
menunjukkan persona laki-laki, sedangkan baris bawah menunjukkan persona
perempuan. Satu titik mewakili hasil satu penanda linguistik pada satu model,
warna titik membedakan kelompok penanda linguistik pada Tabel~\ref{tab:aturan-penanda}, dan
ukuran titik mengikuti jumlah pasangan respons yang memiliki nilai berbeda.
Sumbu horizontal menunjukkan \textit{effect size} \textit{rank-biserial},
dengan nilai positif untuk penanda linguistik yang lebih tinggi pada skenario
persona dan nilai negatif untuk penanda linguistik yang lebih tinggi pada N.
Sumbu vertikal menunjukkan $-\log_{10}(p_{\text{Holm}})$, garis putus-putus
menandai batas signifikansi setelah koreksi Holm, dan nilai di atas 50
ditempatkan pada 50.

\begin{figure}[H]
    \centering
    \includegraphics[width=\textwidth]{image/BabVI/gsm8k/gender/consistent_20261002/03_numeric_effect_volcano_vs_neutral.png}
    \caption[Ukuran efek penanda linguistik pada domain gender]
    {Ukuran efek penanda linguistik numerik terhadap skenario netral dalam
    domain gender}
    \label{fig:marker-gender-terhadap-netral}
\end{figure}

Pada DeepSeek R1-Distill 32B, I1, I2, dan I3 lebih sedikit menggunakan langkah
eksplisit dan LaTeX \textit{display} pada kedua persona. Daftar berpoin dan
operasi aritmetika di dalam kalimat lebih sering digunakan. Pada Qwen3 Next
80B-A3B, ketiga skenario implisit menghasilkan lebih sedikit kata, langkah
eksplisit, subjudul, dan LaTeX \textit{display} pada kedua persona.

Pada Llama 3.3 70B, kata penghubung sebab-akibat dan kata ganti inklusif lebih
sering digunakan pada kedua persona, sedangkan daftar berpoin berkurang. Pada
GPT-OSS 20B, E menghasilkan lebih sedikit kata dan LaTeX \textit{display}
daripada N pada kedua persona, disertai kenaikan variasi kata dan penggunaan
subjudul. Pada I1, I2, dan I3, penggunaan subjudul meningkat,
sedangkan rata-rata panjang kalimat menurun.

Pada keenam model, skenario persona memengaruhi cara model menyusun
penyelesaian. Kesamaan perubahan pada persona laki-laki dan perempuan,
terutama pada DeepSeek R1-Distill 32B dan Qwen3 Next 80B-A3B, menunjukkan bahwa
susunan skenario eksplisit dan implisit ikut menentukan perubahan terhadap N.
Perubahan khusus antara persona laki-laki dan perempuan diperiksa melalui
perbandingan antarpersona.

\par\noindent\textbf{b. Perbandingan Persona Laki-Laki dan Perempuan}\par

Perbandingan antarpersona mencakup 384 hasil, yaitu 16 penanda linguistik pada
empat skenario dan enam model. Setelah koreksi Holm, 56 perbandingan tetap
signifikan dan seluruh model memiliki setidaknya satu hasil signifikan.
Tabel~\ref{tab:penanda-gender-antarpersona} merangkum jumlah hasil signifikan,
skenario, dan temuan dengan $|r|$ terbesar pada setiap model.

\begin{table}[H]
    \centering
    \caption{Ringkasan perbandingan penanda linguistik antara persona
    laki-laki dan perempuan}
    \label{tab:penanda-gender-antarpersona}
    \small
    \renewcommand{\arraystretch}{1.18}
    \setlength{\tabcolsep}{4pt}
    \begin{tabularx}{\textwidth}{
        >{\raggedright\arraybackslash}p{2.65cm}
        >{\centering\arraybackslash}p{1.35cm}
        >{\centering\arraybackslash}p{1.45cm}
        >{\raggedright\arraybackslash}X}
        \toprule
        \textbf{Model} &
        \textbf{Jumlah signifikan} &
        \textbf{Skenario} &
        \textbf{Temuan dengan $|r|$ terbesar} \\
        \midrule
        DeepSeek R1-Distill 32B & 5 & I2, I3 &
        Subjudul lebih sering digunakan oleh persona perempuan pada I2
        ($r=-0{,}18$). \\
        Llama 3.1 8B & 12 & E, I1--I3 &
        Subjudul lebih sering digunakan oleh persona laki-laki pada I3
        ($r=0{,}84$). \\
        Llama 3.3 70B & 12 & E, I1--I3 &
        Subjudul lebih sering digunakan oleh persona laki-laki pada I3
        ($r=0{,}98$). \\
        Ministral 8B & 14 & E, I1--I3 &
        Jumlah kata lebih tinggi pada persona perempuan di I3
        ($r=-0{,}48$). \\
        GPT-OSS 20B & 8 & E, I1--I3 &
        Subjudul lebih sering digunakan oleh persona perempuan pada I3
        ($r=-0{,}35$). \\
        Qwen3 Next 80B-A3B & 5 & I1--I3 &
        Subjudul lebih sering digunakan oleh persona perempuan pada I1
        ($r=-0{,}41$). \\
        \bottomrule
    \end{tabularx}

    \vspace{0.15cm}
    \begin{minipage}{0.98\textwidth}
        \small
        Jumlah signifikan menunjukkan banyaknya kombinasi skenario dan
        penanda linguistik dengan $p_{\mathrm{Holm}}<0{,}05$. Nilai $r$ positif
        menunjukkan penanda linguistik lebih tinggi pada persona laki-laki,
        sedangkan nilai negatif menunjukkan penanda linguistik lebih tinggi
        pada persona perempuan.
    \end{minipage}
\end{table}

Penggunaan subjudul berbeda antara persona laki-laki dan perempuan menurut
model dan skenario. Pada I3, subjudul lebih sering digunakan oleh persona
laki-laki pada kedua model Llama, sedangkan persona perempuan lebih sering
menggunakannya pada GPT-OSS 20B. Persona perempuan juga lebih sering
menggunakan subjudul pada I1 Qwen3 Next 80B-A3B.

Nilai $r$ yang besar pada kedua model Llama di I3 berasal dari sedikit pasangan
yang memiliki jumlah subjudul berbeda. Pada Llama 3.3 70B, 65 dari 68 pasangan
yang berbeda lebih tinggi pada persona laki-laki, sedangkan tiga pasangan lebih
tinggi pada persona perempuan. Pada Llama 3.1 8B, 57 dari 65 pasangan yang
berbeda lebih tinggi pada persona laki-laki, sedangkan delapan pasangan lebih
tinggi pada persona perempuan. Nilai $r$ yang besar menunjukkan bahwa hampir
seluruh pasangan yang berubah memiliki jumlah subjudul lebih tinggi pada persona
laki-laki, meskipun perbedaannya hanya muncul pada sebagian kecil soal.

Pasangan respons Llama 3.3 70B pada soal
\texttt{gsm8k\_01035} memperlihatkan perubahan tersebut. Respons persona
laki-laki memisahkan perhitungan dengan subjudul \textit{Cost of books:},
\textit{Cost of pencils:}, dan \textit{Total amount spent:}, sedangkan respons
persona perempuan menulis \textit{The cost of 3 books at \$16 each is:} dan
\textit{The cost of 3 pencils at \$6 each is:} tanpa subjudul. Kedua respons
menghasilkan total \$66, sehingga perubahan terdapat pada susunan penyelesaian,
bukan pada hasil hitung untuk soal tersebut.

Pada Ministral 8B, I2 menghasilkan respons persona
laki-laki yang lebih panjang dengan variasi kata lebih rendah, sedangkan I3
menghasilkan respons yang lebih pendek dengan variasi kata lebih tinggi.
Kedua skenario berbeda dalam isi percakapan \textit{warmup} yang digunakan.
Pengujian ini belum memisahkan pengaruh identitas gender dari isi
\textit{warmup} terhadap bentuk respons.
Rincian seluruh perbandingan terdapat pada
Lampiran~\ref{fig:lampiran-marker-gender-laki-perempuan}.

\par\noindent\textbf{c. Bentuk Respons dan Bahasa Matematis}\par

Penanda linguistik bentuk respons mencakup jumlah kata, panjang kalimat, variasi kata,
langkah eksplisit, daftar berpoin, dan subjudul. Bahasa matematis pada bagian
ini menunjukkan cara model menempatkan perhitungan dalam respons. LaTeX
\textit{display} menempatkan persamaan pada baris tersendiri, sedangkan operasi
aritmetika di dalam kalimat menyatukan perhitungan seperti \(12+8=20\) dengan
penjelasan. Kedua bentuk tersebut menunjukkan apakah model lebih sering
memisahkan persamaan dari kalimat atau menuliskan perhitungan sebagai bagian
dari kalimat.

DeepSeek R1-Distill 32B menjadi contoh utama pada perbandingan dengan N. I1,
I2, dan I3 lebih sedikit menggunakan langkah eksplisit dan LaTeX
\textit{display}, tetapi lebih sering menggunakan daftar berpoin dan operasi
aritmetika di dalam kalimat. Penyelesaian pada skenario implisit lebih sering
ditulis melalui kalimat dan operasi di dalam teks
daripada melalui urutan langkah dan persamaan terpisah. Pada I2 dan I3,
persona laki-laki juga menggunakan kalimat yang lebih panjang, sedangkan
persona perempuan lebih sering menggunakan daftar berpoin.

Perubahan tersebut terlihat pada soal \texttt{gsm8k\_00238}. Pada N, DeepSeek
memakai tiga label langkah dan menempatkan perhitungan \(20+10=30\) serta
\(30\times3=90\) pada baris tersendiri. Pada I2 persona laki-laki, perhitungan
yang sama diringkas menjadi dua poin, yaitu \textit{Total amount before
investment: \$20 + \$10 = \$30} dan \textit{Amount after tripling: \$30
$\times$ 3 = \$90}.

Pada Ministral 8B, I2 menghasilkan respons persona laki-laki yang lebih
panjang dengan variasi kata yang lebih rendah. Pada I3, respons persona
laki-laki menjadi lebih pendek dengan variasi kata yang lebih tinggi. I2
menyampaikan persona melalui aktivitas merapikan penampilan, sedangkan I3
menggunakan pilihan pakaian. Panjang dan susunan respons berbeda antara
I2 dan I3.

Jumlah kata dan \textit{completion token} mengukur bagian respons yang
berbeda. Pada DeepSeek R1-Distill 32B, teks yang terlihat menggunakan lebih
sedikit kata pada skenario implisit, sementara jumlah \textit{completion
token} meningkat. Jumlah kata dihitung dari teks respons, sedangkan
\textit{completion token} diperoleh dari metadata API. Skenario persona
karena itu dapat mengubah penggunaan token dan bentuk teks yang terlihat
dengan cara yang berbeda.

\par\noindent\textbf{d. Pola Pembuka dan Penanda Linguistik Identitas}\par

Pola pembuka respons dibandingkan pada N, E, I1, I2, dan I3 untuk setiap
model dan persona. Contoh pembuka instruksional pada respons asli adalah
\textit{Let's solve the problem.} pada \texttt{gsm8k\_00384}, sedangkan
pembuka berupa subjudul muncul sebagai \textit{Step-by-Step Solution:} pada
\texttt{gsm8k\_00241}. Pembuka naratif langsung memulai isi soal, seperti
\textit{Emil is 19 now.} pada \texttt{gsm8k\_00980}, sedangkan pengulangan
soal muncul sebagai \textit{The problem states that three-fourths of the
students have a desktop computer} pada \texttt{gsm8k\_00475}. Persentase
setiap kategori ditampilkan pada Gambar~\ref{fig:pembuka-gender-gsm8k}.

\begin{figure}[H]
    \centering
    \includegraphics[width=\textwidth]{image/BabVI/gsm8k/gender/consistent_20261002/08_opening_composition_100pct.png}
    \caption{Persentase kategori pembuka respons pada domain gender}
    \label{fig:pembuka-gender-gsm8k}
\end{figure}

Pada DeepSeek R1-Distill 32B, sebanyak 61,03\% respons pada N menggunakan
pembuka instruksional dan 38,74\% menggunakan pembuka berupa subjudul. Pada E,
pembuka berupa subjudul mencakup 55,74--56,40\% respons, sedangkan pembuka
instruksional mencakup 40,28--41,38\%. Pada I1, I2, dan I3, pembuka naratif
mencakup 54,21--66,49\% respons. Pengulangan soal juga ditemukan pada
18,38--21,87\% respons.

GPT-OSS 20B tetap lebih sering menggunakan pembuka naratif. Persentasenya
mencapai 80,91\% pada N dan 61,32--69,95\% pada I1, I2, dan I3. Pada Qwen3
Next 80B-A3B, pembuka instruksional mencakup 47,16--64,72\% respons pada
kelima skenario.

Pada N, pembuka instruksional digunakan pada 99,32\% respons Llama 3.1 8B,
95,68\% respons Llama 3.3 70B, dan 96,59\% respons Ministral 8B. Sebagian
besar respons ketiga model tersebut sudah menggunakan pembuka instruksional
sebelum persona diberikan. Pembuka ini tetap menjadi kategori utama pada E,
I1, I2, dan I3 dengan persentase 95,14--99,62\%.

Setelah koreksi Holm, lima dari enam model menghasilkan perbedaan komposisi
pembuka antara persona laki-laki dan perempuan. Perbedaan ditemukan pada I2
dan I3 untuk DeepSeek R1-Distill 32B, I3 untuk Llama 3.3 70B dan Ministral 8B,
E, I1, dan I3 untuk GPT-OSS 20B, serta I1 dan I2 untuk Qwen3 Next 80B-A3B.
Llama 3.1 8B tidak menghasilkan perbedaan signifikan pada keempat skenario.
Pada lima model tersebut, persona laki-laki dan perempuan membuat model
memulai penyelesaian dengan bentuk yang berbeda. Rincian setiap perbandingan
dicantumkan pada Lampiran~\ref{fig:lampiran-pembuka-langsung-gender}.

Penanda linguistik identitas digunakan untuk melihat apakah identitas disebut secara
langsung di dalam respons. Pernyataan identitas AI mencakup ungkapan seperti
\textit{As an AI language model}, sedangkan ungkapan pengarah pembaca mencakup
\textit{note that}, \textit{you can see that}, dan \textit{consider this}.
Pronomina gender mencakup \textit{he}, \textit{him}, \textit{his},
\textit{she}, \textit{her}, dan \textit{hers}.

Tidak ada model yang menyatakan dirinya sebagai AI dalam respons GSM8K.
Ungkapan pengarah pembaca hanya ditemukan pada sedikit respons dan tidak
berbeda secara signifikan antarskenario. Pronomina gender ditemukan pada
beberapa model dan skenario, tetapi digunakan untuk menyebut tokoh dalam soal.
Pengaruh persona gender lebih sering muncul pada bentuk
pembuka daripada melalui penyebutan identitas secara langsung.

\par\noindent\textbf{e. Penanda Linguistik dan Akurasi}\par

Pada domain gender, penanda linguistik berubah pada keenam model. E tidak
menghasilkan perubahan akurasi yang signifikan terhadap N, sedangkan perubahan
akurasi ditemukan pada I1 dan I2 pada DeepSeek R1-Distill 32B serta I3 persona
perempuan pada Llama 3.3 70B. Perubahan penanda linguistik tidak selalu diikuti
kenaikan atau penurunan akurasi yang sama.

I1 menyampaikan persona laki-laki sebagai ayah yang melindungi dan menafkahi
keluarga, sedangkan persona perempuan sebagai ibu yang mengatur jadwal,
menyiapkan makanan, dan memperhatikan kebutuhan keluarga. Percakapan
\textit{warmup} I1 paling panjang di antara ketiga skenario implisit. Pada
DeepSeek, I1 lebih sering menghasilkan pembuka naratif atau pengulangan
informasi soal daripada E. Langkah eksplisit dan LaTeX \textit{display}
berkurang, sedangkan daftar berpoin dan operasi aritmetika di dalam kalimat
meningkat. Perubahan bentuk respons tersebut muncul bersama penurunan akurasi
sebesar 2,96 poin untuk persona laki-laki dan 3,18 poin untuk persona perempuan
dibandingkan dengan E.

I2 menyampaikan persona melalui kegiatan bercukur dan merapikan janggut untuk
laki-laki serta merias wajah untuk perempuan. Pada DeepSeek R1-Distill 32B,
I2 juga lebih sering menghasilkan pembuka naratif dan pengulangan soal serta
lebih jarang menggunakan langkah eksplisit dan LaTeX \textit{display}.
Akurasi persona laki-laki lebih rendah 2,28 poin daripada N, sedangkan akurasi
persona perempuan lebih rendah 2,89 poin daripada N dan 2,35 poin daripada E.
Meskipun isi I1 dan I2 berbeda, keduanya memperlihatkan perubahan susunan
respons dan penurunan akurasi pada DeepSeek.

I3 menyampaikan persona melalui jas dan dasi untuk laki-laki serta gaun dan
sepatu hak tinggi untuk perempuan. Pada Llama 3.3 70B, akurasi I3 persona
perempuan lebih tinggi 2,12 poin daripada E. Pada perbandingan yang sama, kata
penghubung sebab-akibat dan kata ganti inklusif meningkat, sedangkan penggunaan
daftar berpoin berkurang. Hasil I3 berbeda dari I1 dan I2 karena perubahan
penanda linguistiknya muncul bersama kenaikan, bukan penurunan, akurasi.
Perbandingan akurasi I3 lainnya tidak signifikan setelah koreksi Holm.

Pada perbandingan antarpersona laki-laki dan perempuan, tidak ada
perbedaan akurasi yang signifikan pada semua model dan skenario. Perbedaannya
lebih sering muncul pada cara jawaban ditulis. Pada I3, kedua model Llama lebih
sering menggunakan subjudul untuk persona laki-laki, sedangkan GPT-OSS 20B
lebih sering menggunakannya untuk persona perempuan. DeepSeek pada I2 dan Qwen
pada I1 juga lebih sering menggunakan subjudul untuk persona perempuan,
sementara pada Ministral I3 jawaban persona perempuan lebih panjang. Pada
GSM8K, perbedaan antara persona laki-laki dan perempuan lebih banyak terdapat
pada cara penyampaian jawaban daripada akurasinya.

\subsubsection{Ringkasan Domain Gender}
\label{subsubsec:ringkasan-gender-gsm8k}

Untuk rumusan masalah pertama, persona gender lebih sering mengubah penanda
linguistik daripada akurasi pada GSM8K. E tidak menghasilkan perbedaan akurasi
yang signifikan terhadap N pada keenam model. I1 dan I2 menurunkan akurasi
kedua persona pada DeepSeek R1-Distill 32B dibandingkan dengan N, sedangkan I3
meningkatkan akurasi persona perempuan pada Llama 3.3 70B dibandingkan dengan
E. Penanda
linguistik berubah pada keenam model meskipun perbandingan akurasi antarpersona
laki-laki dan perempuan tidak signifikan. Persona yang
akurasinya berubah karena itu berbeda menurut cara penyampaian dan model.

I1 menyampaikan persona melalui peran ayah dan ibu, sedangkan I2 menggunakan
kegiatan merapikan penampilan. Kedua skenario tersebut membuat DeepSeek lebih
sering memakai pembuka naratif dan lebih jarang memakai langkah eksplisit
serta LaTeX \textit{display}, disertai penurunan akurasi. I3 menggunakan
pakaian dan pada Llama 3.3 70B persona perempuan menghasilkan lebih banyak
kata penghubung sebab-akibat serta kenaikan akurasi. Setiap skenario implisit
memuat tema persona dan riwayat percakapan, sehingga hasilnya merujuk pada
skenario persona lengkap yang diterima model.

Dari sisi penggunaan sumber daya, skenario implisit menambah \textit{input
token} karena percakapan \textit{warmup} dikirim bersama soal. Pada rata-rata
I1--I3, persona laki-laki menggunakan total token 0,6--26,7\% lebih banyak
daripada persona perempuan pada keenam model. Penggunaan token yang lebih
besar tidak menghasilkan keunggulan akurasi persona laki-laki. Akurasi
rata-rata persona perempuan lebih tinggi 0,37--0,44 poin pada lima dari enam
model, tetapi seluruh perbandingan antarpersona tidak signifikan.

Untuk rumusan masalah kedua, bagian ini menunjukkan hasil domain gender pada
tugas penalaran numerik. Pada domain ini, persona memengaruhi cara penulisan
respons pada lebih banyak model daripada akurasinya. Perbandingan dengan
domain agama dan regional-kultural terdapat pada
Subbab~\ref{subsec:perbandingan-domain-gsm8k}, sedangkan perbandingan dengan
tugas penalaran moral menggunakan hasil ETHICS pada
Subbab~\ref{subsubsec:ringkasan-gender-ethics}.

Untuk rumusan masalah ketiga, pengaruh persona gender tidak seragam pada
berbagai jenis model dan ukuran parameter. Di antara model
\textit{reasoning}, perubahan akurasi terdapat pada DeepSeek, tetapi tidak pada
GPT-OSS 20B dan Qwen3 Next 80B-A3B. Di antara model \textit{non-reasoning},
perubahan akurasi hanya terdapat pada Llama 3.3 70B, sedangkan Llama 3.1 8B
dan Ministral 8B tetap mengubah penanda linguistik. Jenis model dan jumlah
parameter karena itu tidak menghasilkan pengaruh persona yang sama.


% =====================================================================
\subsection{Domain Agama}
\label{subsec:hasil-agama-gsm8k}

Domain agama mencakup persona Islam, Kristen, Yahudi, dan ateis. Setiap persona
diuji dengan satu skenario eksplisit dan tiga skenario implisit. I1 menggunakan
kitab agama sebagai pedoman hidup, I2 menggunakan hari raya keagamaan, dan I3
menggunakan kegiatan ibadah atau komunitas agama; pada persona ateis, ketiga
tema tersebut dibuat dalam bentuk nonkeagamaan. Ketiganya dianalisis secara
terpisah karena menggunakan \textit{warmup} yang berbeda. Analisis menggunakan
134.538 respons dari 1.319 soal GSM8K.

\subsubsection{Akurasi Respons}
\label{subsubsec:akurasi-agama-gsm8k}

\par\noindent\textbf{a. Akurasi Setiap Skenario}\par

Tabel~\ref{tab:akurasi-agama-gsm8k} menunjukkan akurasi keenam model pada
skenario netral serta skenario yang menggunakan persona Islam, Kristen,
Yahudi, dan ateis. Baris N berisi hasil pada skenario netral. Untuk setiap
persona, hasil skenario eksplisit ditulis pada baris E, sedangkan hasil tiga
skenario implisit ditulis pada baris I1, I2, dan I3.

\begin{table}[H]
    \centering
    \caption{Akurasi GSM8K pada domain agama menurut persona dan skenario (\%)}
    \label{tab:akurasi-agama-gsm8k}
    \scriptsize
    \setlength{\tabcolsep}{3pt}
    \renewcommand{\arraystretch}{1.08}
    \begin{tabularx}{\textwidth}{
        >{\raggedright\arraybackslash}p{1.75cm}
        >{\centering\arraybackslash}p{0.65cm}
        *{6}{>{\centering\arraybackslash}X}}
        \toprule
        \textbf{Persona} & \textbf{Skenario} & \textbf{Qwen3} &
        \textbf{Llama 3.1} & \textbf{Ministral} & \textbf{DeepSeek} &
        \textbf{GPT-OSS} & \textbf{Llama 3.3} \\
        \midrule
        Netral & N & 94,08 & 84,43 & 93,47 & 93,93 & 93,93 & 91,96 \\
        \addlinespace[2pt]
        Islam & E  & 96,21 & 81,96 & 91,43 & 92,87 & 90,49 & 93,93 \\
              & I1 & 95,07 & 80,89 & 82,64 & 92,80 & 92,55 & 93,56 \\
              & I2 & 95,45 & 81,20 & 86,73 & 93,10 & 93,24 & 94,01 \\
              & I3 & 95,68 & 79,53 & 84,38 & 93,18 & 92,91 & 94,47 \\
        \addlinespace[2pt]
        Kristen & E  & 96,36 & 84,23 & 91,05 & 92,95 & 91,81 & 94,01 \\
                & I1 & 94,62 & 80,74 & 85,97 & 93,93 & 92,94 & 93,86 \\
                & I2 & 95,60 & 79,76 & 88,48 & 92,42 & 93,62 & 94,09 \\
                & I3 & 95,15 & 83,40 & 86,66 & 92,57 & 92,34 & 94,31 \\
        \addlinespace[2pt]
        Yahudi & E  & 95,53 & 83,47 & 90,60 & 93,25 & 92,23 & 94,01 \\
               & I1 & 95,60 & 83,24 & 87,34 & 92,80 & 92,33 & 93,63 \\
               & I2 & 95,68 & 82,11 & 87,95 & 93,78 & 93,01 & 94,16 \\
               & I3 & 95,44 & 80,89 & 87,72 & 93,48 & 92,86 & 93,56 \\
        \addlinespace[2pt]
        Ateis & E  & 95,68 & 84,00 & 90,22 & 92,65 & 92,09 & 93,93 \\
              & I1 & 95,60 & 82,11 & 87,11 & 92,87 & 90,82 & 93,18 \\
              & I2 & 95,91 & 81,35 & 86,50 & 92,65 & 91,79 & 94,01 \\
              & I3 & 95,75 & 81,88 & 86,13 & 93,71 & 92,34 & 94,84 \\
        \bottomrule
    \end{tabularx}
\end{table}

Dalam satu model, selisih antara akurasi tertinggi dan terendah berkisar
1,51--10,83 poin persentase. Selisih terbesar terdapat pada Ministral 8B,
sedangkan selisih terkecil terdapat pada DeepSeek R1-Distill 32B.
Gambar~\ref{fig:ringkasan-akurasi-agama-gsm8k} menunjukkan perbandingan hasil
pada skenario netral, skenario eksplisit, dan rata-rata I1, I2, dan I3.

\begin{figure}[H]
    \centering
    \includegraphics[
        width=\textwidth,
        height=0.46\textheight,
        keepaspectratio
    ]{image/bab6/gsm8k/religion/akurasi/05_religion_absolute_accuracy_explicit_implicit.png}
    \caption[Akurasi GSM8K menurut skenario pada domain agama]
    {Perbandingan akurasi GSM8K antara skenario netral, skenario eksplisit,
    dan rata-rata skenario implisit pada domain agama}
    \label{fig:ringkasan-akurasi-agama-gsm8k}
\end{figure}

Perbandingan dengan skenario netral membentuk tiga pola. Seluruh skenario
persona lebih tinggi pada Qwen3 Next 80B-A3B dan Llama 3.3 70B, tetapi lebih
rendah pada Llama 3.1 8B, Ministral 8B, dan GPT-OSS 20B. Pada DeepSeek
R1-Distill 32B, seluruh skenario persona lebih rendah, kecuali I1 persona
Kristen yang memperoleh akurasi sama dengan skenario netral. Penurunan terbesar
terdapat pada I1 persona Islam pada Ministral 8B, yaitu 10,83 poin.

Gambar~\ref{fig:delta-implisit-eksplisit-agama-gsm8k} membandingkan rata-rata
I1--I3 dengan E untuk setiap persona. Nilai positif menunjukkan rata-rata
akurasi I1--I3 lebih tinggi daripada E, sedangkan nilai negatif menunjukkan
hasil yang lebih rendah.

\begin{figure}[H]
    \centering
    \includegraphics[
        width=\textwidth,
        height=0.46\textheight,
        keepaspectratio
    ]{image/bab6/gsm8k/religion/akurasi/06_religion_delta_implicit_vs_explicit.png}
    \caption[Selisih akurasi antarskenario pada domain agama]
    {Selisih rata-rata akurasi antara skenario implisit dan eksplisit pada
    domain agama}
    \label{fig:delta-implisit-eksplisit-agama-gsm8k}
\end{figure}

Perubahan rata-rata I1--I3 terhadap E paling besar ditemukan pada Ministral
8B. Pada keempat persona, rata-rata akurasinya lebih rendah 2,93--6,85 poin
daripada E. Penurunan juga ditemukan pada Llama 3.1 8B sebesar 1,39--2,93
poin. Sebaliknya, pada GPT-OSS 20B, rata-rata akurasi I1--I3 lebih tinggi
0,50--2,41 poin untuk persona Islam, Kristen, dan Yahudi, tetapi lebih rendah
0,44 poin untuk persona ateis. Pada Qwen3 Next 80B-A3B, DeepSeek R1-Distill
32B, dan Llama 3.3 70B, selisihnya tidak melebihi 1,24 poin. Artinya, skenario
implisit dapat menghasilkan kenaikan atau penurunan akurasi menurut model dan
persona yang digunakan.

\par\noindent\textbf{b. Hasil Perbandingan Statistik}\par
\label{subsubsec:perbandingan-agama-gsm8k}

Pengujian utama membandingkan akurasi E dan I1--I3 pada model dan persona agama
yang sama. Q Cochran digunakan untuk memeriksa perbedaan keempat skenario
tersebut. Jika hasilnya signifikan, I1--I3 dibandingkan dengan E menggunakan
Exact McNemar dan koreksi Holm. Pengujian utama mencakup 72 perbandingan I--E.
Sebanyak 96 perbandingan skenario persona dengan N ditempatkan sebagai analisis
sensitivitas, sedangkan 144 perbandingan antarpersona digunakan untuk
membandingkan persona agama. Kedua kelompok tersebut juga diuji menggunakan Exact
McNemar dan koreksi Holm.

Uji Q Cochran signifikan pada delapan dari 24 kelompok. Empat kelompok terdapat
pada Ministral 8B dan mencakup seluruh persona agama. Empat kelompok lainnya
terdapat pada persona ateis Llama 3.3 70B, persona Kristen Qwen3 Next 80B-A3B
dan Llama 3.1 8B, serta persona Islam GPT-OSS 20B. Setelah penyaringan berdasarkan
Q Cochran dan koreksi Holm, 19 perbandingan I--E tetap signifikan. Pada analisis
sensitivitas, 23 perbandingan dengan N tetap signifikan setelah koreksi Holm.
Empat perbandingan antarpersona juga tetap signifikan.

Dari 96 perbandingan dengan N, 23 tetap signifikan setelah koreksi Holm.
Sebanyak 15 hasil terdapat pada Ministral 8B. Delapan hasil lainnya tersebar
pada empat model. DeepSeek R1-Distill 32B tidak menghasilkan perbandingan
signifikan terhadap N.

Empat dari 144 perbandingan antarpersona tetap signifikan dan seluruhnya
menunjukkan akurasi persona Islam yang lebih rendah. Pada Ministral 8B, akurasi
persona Islam lebih rendah 4,47 poin daripada ateis pada I1 serta lebih rendah
4,70 dan 3,34 poin daripada persona Yahudi pada I1 dan I3. Pada Llama 3.1 8B,
akurasi persona Islam pada I3 lebih rendah 3,87 poin daripada persona Kristen.

Gambar~\ref{fig:effect-agama-gsm8k} menunjukkan \textit{effect size} Cohen's
$h$ pada 52 jenis perbandingan per model dalam empat panel. Panel pertama dan
kedua memuat perbandingan skenario untuk setiap persona, sedangkan panel ketiga
dan keempat memuat perbandingan antarpersona. Warna biru menunjukkan akurasi
skenario atau persona pertama lebih rendah dan warna merah menunjukkan akurasi
yang lebih tinggi. Tanda bintang menunjukkan hasil yang tetap signifikan
setelah koreksi Holm, tanda $\dagger$ menunjukkan perbandingan dengan N, dan
tanda $\ddagger$ menunjukkan perbandingan antarpersona; nilai
\textit{p-value} lengkap terdapat pada
Lampiran~\ref{chap:lampiran-statistik-agama}.

\begin{figure}[H]
    \centering
    \includegraphics[
        width=\linewidth,
        height=0.84\textheight,
        keepaspectratio,
        trim=0 90bp 0 0,
        clip
    ]{image/bab6/gsm8k/religion/akurasi/26_religion_all_52_effect_holm_stacked_26_26.png}
    \caption[Ukuran efek dan signifikansi akurasi pada domain agama]
    {Ukuran efek Cohen's $h$ dan hasil uji Exact McNemar setelah koreksi Holm
    pada 52 jenis perbandingan per model dalam domain agama.\\[-2pt]
    {\footnotesize\textit{Catatan:} * = signifikan setelah koreksi Holm
    ($p_{\mathrm{Holm}} < 0{,}05$); $\dagger$ = perbandingan terhadap
    skenario netral bersama (analisis sensitivitas; protokol berbeda);
    $\ddagger$ = perbandingan antaridentitas. E = eksplisit; I1--I3 = implisit;
    N = netral. Ath = ateis; Chr = Kristen; Isl = Islam; Jew = Yahudi.
    Cohen's $h$ dihitung dari skenario pertama dikurangi skenario pembanding.}}
    \label{fig:effect-agama-gsm8k}
\end{figure}

Tabel~\ref{tab:hasil-signifikan-agama-gsm8k} merangkum 19 perbandingan I--E
dari pengujian utama yang didukung Q Cochran dan tetap signifikan setelah
koreksi Holm. Hasil I1, I2, dan I3 pada model dan persona yang sama digabungkan
dalam satu baris jika seluruh akurasinya naik atau seluruhnya turun. Kolom
selisih dan Cohen's $h$ menunjukkan rentang hasil dari skenario yang tercantum.

\begin{table}[H]
    \centering
    \caption[Perbandingan akurasi I--E pada domain agama]
    {Perbandingan akurasi I--E pada domain agama yang didukung Q Cochran dan
    tetap signifikan setelah koreksi Holm}
    \label{tab:hasil-signifikan-agama-gsm8k}
    \scriptsize
    \renewcommand{\arraystretch}{1.2}
    \begin{adjustbox}{max width=\textwidth}
    \begin{tabular}{>{\raggedright\arraybackslash}p{1.8cm}
                    >{\raggedright\arraybackslash}p{4.2cm}
                    >{\centering\arraybackslash}p{1.3cm}
                    >{\centering\arraybackslash}p{3.0cm}
                    >{\centering\arraybackslash}p{2.8cm}}
        \toprule
        \textbf{Persona} & \textbf{Model} & \textbf{Skenario} &
        \textbf{Selisih I--E} & \textbf{Rentang $h$} \\
        \midrule
        Ateis & Ministral 8B & I1--I3 &
        $-4{,}09$ s.d. $-3{,}11$ & $-0{,}13$ s.d. $-0{,}10$ \\
        Kristen & Llama 3.1 8B & I1--I2 &
        $-4{,}47$ s.d. $-3{,}49$ & $-0{,}12$ s.d. $-0{,}09$ \\
        Kristen & Ministral 8B & I1--I3 &
        $-5{,}08$ s.d. $-2{,}58$ & $-0{,}16$ s.d. $-0{,}09$ \\
        Kristen & Qwen3 Next 80B-A3B & I1, I3 &
        $-1{,}74$ s.d. $-1{,}21$ & $-0{,}08$ s.d. $-0{,}06$ \\
        Islam & Ministral 8B & I1--I3 &
        $-8{,}79$ s.d. $-4{,}70$ & $-0{,}27$ s.d. $-0{,}15$ \\
        Islam & GPT-OSS 20B & I1--I3 &
        $+2{,}06$ s.d. $+2{,}67$ & $+0{,}07$ s.d. $+0{,}10$ \\
        Yahudi & Ministral 8B & I1--I3 &
        $-3{,}26$ s.d. $-2{,}65$ & $-0{,}10$ s.d. $-0{,}09$ \\
        \bottomrule
    \end{tabular}
    \end{adjustbox}
\end{table}

Dari 19 perbandingan I--E yang signifikan, 12 terdapat pada Ministral 8B dan
seluruhnya menunjukkan penurunan 2,58--8,79 poin. Penurunan terbesar terdapat
pada I1 persona Islam sebesar 8,79 poin dengan $h=-0{,}27$, sedangkan nilai
$|h|$ di atas 0,20 hanya terdapat pada I1 dan I3 persona Islam pada Ministral
8B. Tiga hasil GPT-OSS 20B menunjukkan kenaikan 2,06--2,67 poin pada persona
Islam dan menjadi satu-satunya kenaikan pada ketiga skenario implisit untuk
persona yang sama. Empat hasil lainnya terdapat pada persona Kristen pada
Llama 3.1 8B dan Qwen3 Next 80B-A3B dan seluruhnya berupa penurunan.

Perubahan terbesar ditemukan pada persona Islam, tetapi arahnya tidak sama pada
setiap model. Pada Ministral 8B, akurasi I1--I3 persona Islam lebih rendah
daripada E dan tiga perbandingan antarpersona juga menunjukkan akurasi yang lebih
rendah daripada persona ateis dan Yahudi. Pada Llama 3.1 8B, akurasi I3 persona
Islam lebih rendah daripada persona Kristen. Sebaliknya, akurasi I1--I3 persona
Islam pada GPT-OSS 20B meningkat 2,06--2,67 poin dibandingkan dengan E. Artinya,
perubahan akurasi persona Islam tidak sama pada setiap model dan skenario.

\par\noindent\textbf{c. Penggunaan Token}\par
\label{subsubsec:token-agama-gsm8k}

Ketiga skenario implisit memberi konteks agama yang berbeda: I1 menggunakan
sumber pedoman hidup, I2 menggunakan hari raya, dan I3 menggunakan kegiatan
ibadah atau komunitas. Seluruh percakapan \textit{warmup} dikirim bersama soal
sehingga rata-rata \textit{input token} I1--I3 meningkat 436,4--2.309,8\%
dibandingkan dengan E. Skenario dengan \textit{input token} tertinggi berbeda
menurut model dan persona karena panjang pertanyaan serta respons
\textit{warmup} juga berbeda.

\begin{table}[H]
    \centering
    \caption[Perubahan completion token pada domain agama]
    {Rentang perubahan rata-rata \textit{completion token} I1--I3 terhadap E
    pada empat persona agama}
    \label{tab:perubahan-token-agama}
    \small
    \renewcommand{\arraystretch}{1.15}
    \begin{tabular}{lc}
        \toprule
        \textbf{Model} & \textbf{Rentang perubahan} \\
        \midrule
        Qwen3 Next 80B-A3B      & $+9{,}6\%$ sampai $+13{,}7\%$ \\
        Llama 3.1 8B            & $-10{,}8\%$ sampai $+0{,}8\%$ \\
        Ministral 8B            & $-2{,}9\%$ sampai $+13{,}0\%$ \\
        DeepSeek R1-Distill 32B & $+80{,}4\%$ sampai $+86{,}6\%$ \\
        GPT-OSS 20B             & $-19{,}3\%$ sampai $-9{,}9\%$ \\
        Llama 3.3 70B           & $-8{,}9\%$ sampai $+4{,}5\%$ \\
        \bottomrule
    \end{tabular}

    \vspace{0.15cm}
    \begin{minipage}{0.88\textwidth}
        \scriptsize
        Rentang dihitung dari perubahan rata-rata I1--I3 terhadap E pada
        persona Islam, Kristen, Yahudi, dan ateis.
    \end{minipage}
\end{table}

Perubahan \textit{completion token} berbeda antarmodel. Pada keempat persona,
\textit{completion token} meningkat pada DeepSeek R1-Distill 32B dan menurun
pada GPT-OSS 20B. Pada empat model lainnya, kenaikan dan penurunan ditemukan
pada kombinasi model dan persona yang berbeda.
Seluruh model menerima tambahan percakapan \textit{warmup}, tetapi panjang
respons yang dihasilkan dapat meningkat atau menurun.

Jika \textit{input token} dan \textit{completion token} dijumlahkan, rata-rata
total token pada I1--I3 lebih tinggi 129,2--409,5\% daripada E. Kenaikan
tertinggi terdapat pada persona Islam pada Ministral 8B, sedangkan kenaikan
terendah terdapat pada persona Islam pada Qwen3 Next 80B-A3B. Tambahan
\textit{input token} menjadi sumber utama kenaikan total token pada skenario
implisit.

Pada rata-rata I1--I3, selisih total token antara persona dengan jumlah
tertinggi dan terendah dalam satu model berkisar 5,9--43,0\%. Persona Islam
menggunakan total token tertinggi pada tiga model, persona Yahudi pada dua
model, dan persona ateis pada satu model. Selisih terbesar terdapat pada Qwen3
Next 80B-A3B, yaitu 1.019 token antara persona Yahudi dan Islam. Selisih
tersebut terutama berasal dari \textit{input token} karena rata-rata
\textit{completion token} kedua persona hanya berbeda sekitar tiga token.

Penggunaan total token yang lebih besar tidak selalu disertai akurasi yang
lebih tinggi. Pada Ministral 8B dan Llama 3.1 8B, persona Islam menggunakan
total token tertinggi, tetapi memperoleh akurasi lebih rendah pada perbandingan
antarpersona yang signifikan. Pada GPT-OSS 20B, akurasi persona Islam meningkat,
sedangkan total token tertinggi digunakan oleh persona ateis.

Gambar~\ref{fig:lampiran-token-agama-jumlah} menunjukkan jumlah
\textit{completion token} dan \textit{reasoning token} yang tercatat.
Gambar~\ref{fig:lampiran-token-agama-implisit-eksplisit} menunjukkan perubahan
\textit{completion token} I1--I3 terhadap E. Kedua gambar tersebut terdapat
pada Lampiran~\ref{chap:lampiran-statistik-agama}.

\subsubsection{Penanda Linguistik}
\label{subsubsec:penanda-agama-gsm8k}

Perbandingan penanda linguistik menggunakan skenario dan pasangan persona yang
sama dengan pengujian akurasi. Pada GSM8K, perubahan penanda linguistik agama
terdapat pada kosakata dan bentuk respons. E lebih sering memunculkan bahasa
keagamaan, sedangkan I1, I2, dan I3 lebih banyak mengubah panjang respons,
susunan penyelesaian, pembuka, dan bahasa matematis. Penanda linguistik yang berubah
berbeda menurut model dan cara persona disampaikan. Perbandingan antarpersona
juga menghasilkan perubahan yang berbeda, sehingga pengaruh
persona agama muncul melalui kosakata serta cara setiap model menulis respons.

\par\noindent\textbf{a. Perubahan terhadap Skenario Netral}\par

Gambar~\ref{fig:marker-agama-terhadap-netral} dan
Gambar~\ref{fig:marker-agama-terhadap-netral-lanjutan} membandingkan E, I1,
I2, dan I3 dengan N pada persona yang sama. Gambar pertama memuat persona
Islam dan Kristen, sedangkan gambar kedua memuat persona Yahudi dan ateis.
Empat kolom pada setiap gambar menunjukkan E--N, I1--N, I2--N, dan I3--N.
Satu titik mewakili hasil satu penanda linguistik pada satu model, warna titik
membedakan kelompok penanda linguistik pada Tabel~\ref{tab:aturan-penanda},
dan ukuran titik mengikuti jumlah pasangan respons yang memiliki nilai
berbeda. Sumbu horizontal menunjukkan \textit{effect size}
\textit{rank-biserial}; nilai positif menunjukkan penanda linguistik lebih
tinggi pada skenario persona, sedangkan nilai negatif menunjukkan penanda
linguistik lebih tinggi pada N. Sumbu vertikal menunjukkan
$-\log_{10}(p_{\text{Holm}})$, garis putus-putus menandai batas signifikansi
setelah koreksi Holm, dan nilai di atas 50 ditempatkan pada 50.

\begin{figure}[H]
    \centering
    \includegraphics[width=\textwidth]
    {image/BabVI/gsm8k/agama/readable_v3/15_general_numeric_effect_volcano_vs_neutral_islam_christian.pdf}
    \caption[Ukuran efek penanda linguistik pada persona Islam dan Kristen]
    {Ukuran efek penanda linguistik umum terhadap skenario netral dalam
    domain agama pada persona Islam dan Kristen}
    \label{fig:marker-agama-terhadap-netral}
\end{figure}

\begin{figure}[H]
    \centering
    \includegraphics[width=\textwidth]
    {image/BabVI/gsm8k/agama/readable_v3/15_general_numeric_effect_volcano_vs_neutral_jewish_atheist.pdf}
    \caption[Ukuran efek penanda linguistik pada persona Yahudi dan ateis]
    {Ukuran efek penanda linguistik umum terhadap skenario netral dalam
    domain agama pada persona Yahudi dan ateis}
    \label{fig:marker-agama-terhadap-netral-lanjutan}
\end{figure}

Perubahan yang muncul pada keempat persona terdapat pada panjang
kalimat dan susunan penyelesaian. Rata-rata panjang kalimat lebih tinggi dan
daftar berpoin lebih rendah pada seluruh 96 kombinasi empat skenario, empat
persona, dan enam model. Langkah eksplisit juga lebih rendah pada 94 kombinasi.
Dua hasil yang tidak signifikan terdapat pada I1 persona Kristen dan I3
persona ateis pada GPT-OSS 20B.

Pada DeepSeek R1-Distill 32B dan Qwen3 Next 80B-A3B, I1, I2, dan I3 memuat
lebih sedikit kata daripada N pada keempat persona. DeepSeek memakai kalimat
lebih panjang dan variasi kata (TTR-300) lebih tinggi, tetapi lebih jarang
memakai langkah eksplisit, daftar berpoin, subjudul, dan LaTeX
\textit{display}. Operasi aritmetika di dalam kalimat lebih sering digunakan.
Pada Qwen, kata penghubung sebab-akibat, kata ganti orang pertama, keraguan,
dan koreksi diri juga berkurang.

Pada Llama 3.3 70B, kata penghubung sebab-akibat dan kata ganti inklusif lebih
sering digunakan pada I1, I2, dan I3 untuk keempat persona, sedangkan langkah
eksplisit dan daftar berpoin berkurang. Llama 3.1 8B dan Ministral 8B menulis
kalimat yang lebih panjang serta lebih jarang memakai langkah eksplisit dan
daftar berpoin. Pada Ministral, jumlah kata meningkat dan variasi kata
(TTR-300) menurun pada seluruh persona dan skenario implisit.

Kesamaan perubahan pada beberapa identitas, terutama pada DeepSeek
R1-Distill 32B dan Qwen3 Next 80B-A3B, menunjukkan bahwa susunan E dan
percakapan \textit{warmup} ikut membedakan respons dari N. Penanda linguistik yang berubah
tetap tidak sama pada seluruh model. Perubahan khusus antara Islam, Kristen,
Yahudi, dan ateis diperiksa melalui perbandingan antarpersona.

\par\noindent\textbf{b. Perbandingan Antarpersona Agama}\par

Perbandingan antarpersona mencakup Islam--Kristen, Islam--Yahudi,
Islam--ateis, Kristen--Yahudi, Kristen--ateis, dan Yahudi--ateis. Pada 16
penanda linguistik umum, 261 dari 2.304 perbandingan tetap signifikan setelah koreksi
Holm dan seluruh model memiliki setidaknya satu hasil signifikan.
Tabel~\ref{tab:penanda-agama-antaridentitas} merangkum jumlah hasil signifikan,
skenario, dan temuan dengan $|r|$ terbesar pada setiap model.

\begin{table}[H]
    \centering
    \caption{Ringkasan perbandingan penanda linguistik antarpersona agama}
    \label{tab:penanda-agama-antaridentitas}
    \small
    \renewcommand{\arraystretch}{1.18}
    \setlength{\tabcolsep}{4pt}
    \begin{tabularx}{\textwidth}{
        >{\raggedright\arraybackslash}p{2.65cm}
        >{\centering\arraybackslash}p{1.35cm}
        >{\centering\arraybackslash}p{1.45cm}
        >{\raggedright\arraybackslash}X}
        \toprule
        \textbf{Model} &
        \textbf{Jumlah signifikan} &
        \textbf{Skenario} &
        \textbf{Temuan dengan $|r|$ terbesar} \\
        \midrule
        DeepSeek R1-Distill 32B & 25 & E, I1, I3 &
        Penegas lebih sering digunakan oleh persona Kristen daripada Islam
        pada E ($r=-0{,}89$). \\
        Llama 3.1 8B & 57 & E--I3 &
        Subjudul lebih sering digunakan oleh persona ateis daripada Yahudi
        pada I1 ($r=-0{,}995$). \\
        Llama 3.3 70B & 62 & E--I3 &
        Keraguan lebih sering digunakan oleh persona ateis daripada Islam
        pada E ($r=-0{,}999$). \\
        Ministral 8B & 70 & E--I3 &
        Penegas lebih sering digunakan oleh persona Kristen daripada ateis
        pada E ($r=+0{,}97$). \\
        GPT-OSS 20B & 26 & E, I1 &
        Langkah eksplisit lebih sering digunakan oleh persona Kristen daripada
        Islam pada I1 ($r=-1{,}00$). \\
        Qwen3 Next 80B-A3B & 21 & E--I3 &
        Persona Yahudi memakai tiga kata lebih banyak daripada Kristen pada I1
        ($r=-0{,}32$). \\
        \bottomrule
    \end{tabularx}

    \vspace{0.15cm}
    \begin{minipage}{0.98\textwidth}
        \small
        Jumlah signifikan menunjukkan banyaknya kombinasi pasangan identitas,
        skenario, dan penanda linguistik dengan $p_{\mathrm{Holm}}<0{,}05$. Nilai $r$
        positif menunjukkan penanda linguistik lebih tinggi pada persona pertama,
        sedangkan nilai negatif menunjukkan penanda linguistik lebih tinggi pada persona
        kedua.
    \end{minipage}
\end{table}

Perbedaan paling banyak terdapat pada E, yaitu 117 hasil. I1, I2, dan I3
masing-masing memuat 62, 43, dan 39 hasil. Pada E, identitas agama disebutkan
dalam pesan yang sama dengan soal, sedangkan I1--I3 menyampaikan persona lebih
dahulu melalui percakapan \textit{warmup}. Jumlah kata, variasi kata
(TTR-300), rata-rata panjang kalimat, dan penanda linguistik kepastian menjadi pembeda
utama pada hasil tersebut.

Nilai $r$ yang besar perlu dilihat bersama jumlah pasangan yang memiliki nilai
berbeda. Pada I1 Llama 3.1 8B, subjudul hanya berbeda pada 28 dari 1.319
pasangan persona Yahudi dan ateis. Persona ateis memakai lebih banyak subjudul
pada 27 pasangan, sedangkan persona Yahudi lebih tinggi pada satu pasangan.
Sebanyak 1.291 pasangan lainnya memiliki jumlah subjudul yang sama.

Pasangan respons pada soal \texttt{gsm8k\_01009} memperlihatkan perubahan
tersebut. Respons persona Yahudi menuliskan perhitungan jarak Matteo dan Shandy
dalam paragraf tanpa subjudul, sedangkan persona ateis memakai subjudul
\textit{Matteo's distance:} dan \textit{Shandy's distance:}. Kedua respons
memperoleh selisih 230 mil, sehingga perbedaannya terdapat pada susunan
penyelesaian untuk soal tersebut.

Perbedaan antarpersona tidak tetap sama pada I1, I2, dan I3. DeepSeek
R1-Distill 32B tidak memiliki hasil signifikan pada I2, sedangkan GPT-OSS 20B
hanya memilikinya pada E dan I1. Percakapan I1 membahas pedoman hidup, I2
membahas hari raya, dan I3 membahas ibadah atau komunitas. Perubahan yang
berbeda antarskenario menempatkan isi percakapan \textit{warmup} bersama
identitas agama sebagai bagian dari persona yang diterima model. Rincian semua
pasangan terdapat pada Lampiran~\ref{fig:lampiran-penanda-umum-agama-e} sampai
Gambar~\ref{fig:lampiran-penanda-umum-agama-i3}.

\par\noindent\textbf{c. Bentuk Respons dan Bahasa Matematis}\par

Penanda linguistik bentuk respons mencakup jumlah kata, panjang kalimat, variasi kata,
langkah eksplisit, daftar berpoin, dan subjudul. Bahasa matematis menunjukkan
cara model menempatkan perhitungan dalam respons. LaTeX \textit{display}
menempatkan persamaan pada baris tersendiri, sedangkan operasi aritmetika di
dalam kalimat menyatukan perhitungan seperti \(8\times3=24\) dengan
penjelasan. Kedua bentuk tersebut membedakan persamaan yang dipisahkan dari
kalimat dan perhitungan yang ditulis sebagai bagian dari kalimat.

I1 menyampaikan persona melalui pedoman hidup, I2 melalui hari raya, dan I3
melalui ibadah atau komunitas. Perbandingan dengan E menghasilkan 157 hasil
signifikan pada I1, 152 pada I2, dan 158 pada I3 untuk 16 penanda linguistik umum.
Jumlah yang berdekatan menunjukkan bahwa ketiga cara penyampaian sama-sama
disertai perubahan bentuk respons. Penanda linguistik yang meningkat atau menurun tetap
berbeda menurut model dan identitas.

Llama 3.1 8B, Llama 3.3 70B, Ministral 8B, dan Qwen3 Next 80B-A3B lebih sering
membagi respons implisit menjadi langkah, daftar berpoin, atau subjudul
daripada E. Pada DeepSeek R1-Distill 32B, respons implisit lebih singkat,
memakai kalimat lebih pendek, dan lebih sering memakai daftar berpoin daripada
E. GPT-OSS 20B paling sering membedakan respons implisit melalui kalimat yang
lebih panjang.

Bahasa matematis juga berubah menurut model. Pada seluruh 12 perbandingan
DeepSeek, I1, I2, dan I3 lebih jarang menempatkan persamaan pada baris
tersendiri dan lebih sering menulis operasi aritmetika di dalam kalimat.
Qwen lebih jarang memakai kedua bentuk tersebut pada perbandingan yang
signifikan, sedangkan GPT-OSS lebih sering memakai LaTeX \textit{display}.
Llama 3.1 lebih sering menulis operasi aritmetika di dalam kalimat pada lima
perbandingan yang signifikan.

Perubahan DeepSeek terlihat pada soal \texttt{gsm8k\_00009}. Respons N memakai
tiga label langkah dan tujuh persamaan LaTeX \textit{display} untuk menghitung
jarak John dari rumah. Pada I2 persona Islam, perhitungan yang sama ditempatkan
di dalam kalimat setelah subjudul \textit{Step-by-Step Explanation}, termasuk
\(60\times3=180\) dan \(180-135=45\). Kedua respons memperoleh angka 45 mil,
tetapi I2 tidak memisahkan persamaan ke baris tersendiri.

Jumlah kata dan \textit{completion token} mengukur bagian respons yang
berbeda. Pada DeepSeek R1-Distill 32B, teks respons implisit memuat lebih
sedikit kata daripada E, sementara \textit{completion token} meningkat
80,4--86,6\%. Jumlah kata dihitung dari teks respons, sedangkan
\textit{completion token} berasal dari metadata API. Satu token dapat berupa
potongan kata, simbol, atau bagian dari rumus, sehingga kedua ukuran tersebut
dapat berubah dengan besar yang berbeda.

\par\noindent\textbf{d. Pola Pembuka dan Bahasa Keagamaan}\par

Pola pembuka respons dibandingkan pada N, E, I1, I2, dan I3 untuk setiap model
dan persona. Contoh pembuka instruksional pada respons asli adalah
\textit{Let's solve the problem.} pada \texttt{gsm8k\_00384}, sedangkan
pembuka berupa subjudul muncul sebagai \textit{Step-by-Step Solution:} pada
\texttt{gsm8k\_00241}. Pembuka naratif langsung memulai isi soal, seperti
\textit{Emil is 19 now.} pada \texttt{gsm8k\_00980}, sedangkan pengulangan
soal muncul sebagai \textit{The problem states that three-fourths of the
students have a desktop computer} pada \texttt{gsm8k\_00475}. Persentase
setiap kategori ditampilkan pada Gambar~\ref{fig:pembuka-agama-gsm8k}.

\begin{figure}[H]
    \centering
    \includegraphics[
        width=\textwidth,
        height=0.78\textheight,
        keepaspectratio
    ]{image/BabVI/gsm8k/agama/11_opening_composition_100pct.png}
    \caption{Persentase kategori pembuka respons pada domain agama}
    \label{fig:pembuka-agama-gsm8k}
\end{figure}

Pada DeepSeek R1-Distill 32B, sebanyak 61,03\% respons pada N menggunakan
pembuka instruksional dan 38,74\% menggunakan pembuka subjudul. Pada E,
pembuka subjudul mencakup 43,90--56,94\% respons, sedangkan pembuka
naratif mencakup 23,12--39,73\%. Pada I1, I2, dan I3, pembuka naratif
mencakup 66,79--69,87\% respons. Pengulangan soal juga ditemukan pada
20,90--23,53\% respons.

GPT-OSS 20B lebih sering menggunakan pembuka naratif. Persentasenya mencapai
80,91\% pada N dan 62,81--73,40\% pada E, I1, I2, dan I3. Pada Qwen3 Next
80B-A3B, pembuka instruksional mencakup 64,72\% respons pada N dan
42,14--54,34\% pada keempat skenario persona.

Pada Llama 3.1 8B, pembuka instruksional mencakup 99,17--99,77\% respons pada
N dan I1--I3, sedangkan pembuka naratif mencakup 97,50--99,47\% respons pada
E. Pada Llama 3.3 70B, pembuka naratif pada E berkisar 39,27--87,57\% menurut
persona. Pada Ministral 8B, pembuka instruksional pada E berkisar
44,73--99,24\%, sedangkan pembuka naratif berkisar 0,76--47,69\%. Perbedaan
rentang pada E memperlihatkan bahwa identitas yang disebutkan secara langsung
disertai bentuk pembuka yang berbeda.

Setelah koreksi Holm, perbedaan pembuka antarpersona ditemukan pada 30 dari
36 perbandingan E. I1, I2, dan I3 masing-masing memiliki 14, 12, dan 10
perbandingan signifikan. GPT-OSS 20B dan Qwen3 Next 80B-A3B masing-masing
memiliki sembilan hasil pada skenario implisit, sedangkan Llama 3.1 8B tidak
memiliki hasil signifikan pada I1--I3. Rincian setiap pasangan terdapat pada
Lampiran~\ref{fig:lampiran-pembuka-langsung-agama}.

Bahasa keagamaan diukur melalui penyebutan identitas, Tuhan, kitab, kegiatan
keagamaan, ungkapan agama, doa, dan penghormatan terhadap persona. Kata
\textit{Allah}, \textit{God}, \textit{Bible}, \textit{Torah}, dan
\textit{Shabbat} ditemukan dalam respons. Ungkapan \textit{Amen},
\textit{Assalamu Alaikum}, dan \textit{Shalom} juga ditemukan, sehingga istilah
tersebut bukan hanya contoh aturan penanda linguistik. Daftar istilah dan jumlah
kemunculannya terdapat pada Lampiran~\ref{tab:lampiran-ungkapan-agama-audit}.

Perbedaan bahasa keagamaan antarpersona hanya ditemukan pada E. Setelah
koreksi Holm, 115 perbandingan tetap signifikan, yaitu 42 pada Llama 3.1 8B,
35 pada Llama 3.3 70B, 34 pada Ministral 8B, dan empat pada Qwen3 Next
80B-A3B. Tidak ada perbandingan yang signifikan pada DeepSeek R1-Distill 32B
dan GPT-OSS 20B. Pada E, identitas agama disebutkan secara langsung sebelum
soal, sedangkan pada I1--I3 soal utama tidak menyebutkannya.

Istilah pada E juga mengikuti persona yang digunakan. Persona Islam
memunculkan \textit{Allah} dan \textit{Muslim}, persona Kristen memunculkan
\textit{God} dan \textit{Bible}, persona Yahudi memunculkan \textit{Jewish}
dan \textit{Shabbat}, sedangkan persona ateis memunculkan \textit{atheist} dan
\textit{religion}. Tabel istilah serta perbandingan lengkap antarpersona
terdapat pada Lampiran~\ref{tab:lampiran-ungkapan-agama-audit} dan
Gambar~\ref{fig:lampiran-marker-langsung-agama-eksplisit}.

Pada I1 dan I3, istilah agama tidak ditemukan untuk keempat persona. I2 juga
tidak memuat istilah agama pada persona Islam, Kristen, dan ateis. Satu respons
I2 persona Yahudi pada Qwen3 Next 80B-A3B mengulang pembahasan Yom Kippur dari
percakapan \textit{warmup}, sehingga \textit{God}, \textit{Jewish},
\textit{Prayer}, dan \textit{Yom Kippur} masuk ke respons soal. Temuan satu
respons tersebut dicatat terpisah pada
Lampiran~\ref{tab:lampiran-istilah-agama-implisit}.

\par\noindent\textbf{e. Penanda Linguistik dan Akurasi}\par

Persona agama lebih konsisten mengubah bahasa respons daripada akurasi GSM8K
pada model yang diuji. E terutama disertai kosakata dan pembuka yang mengikuti
identitas agama, sedangkan I1, I2, dan I3 disertai perubahan panjang respons,
pembagian langkah, dan bahasa matematis. Perbandingan I--E menghasilkan 19
perubahan akurasi yang signifikan setelah koreksi Holm. Perubahan tersebut
terdapat pada Ministral 8B, GPT-OSS 20B, Llama 3.1 8B, dan Qwen3 Next
80B-A3B.

I1 menyampaikan persona melalui sumber pedoman hidup, seperti Al-Qur'an dan
hadis, Alkitab, Taurat dan penafsiran rabi, serta buku atau penelitian untuk
persona ateis. Pada Ministral 8B, I1 lebih sering memakai langkah eksplisit
dan daftar berpoin serta memiliki variasi kata (TTR-300) yang lebih rendah
daripada E. Perubahan bentuk respons tersebut disertai penurunan akurasi
3,11--8,79 poin pada keempat persona. I1 juga disertai kenaikan 2,06 poin pada
persona Islam GPT-OSS serta penurunan 3,49 dan 1,74 poin pada persona Kristen
Llama 3.1 dan Qwen.

I2 menyampaikan persona melalui Ramadan dan puasa, Natal atau Paskah,
Yom Kippur, serta tidak adanya hari raya agama pada persona ateis. Pada
Ministral 8B, perubahan susunan respons disertai penurunan akurasi
2,58--4,70 poin pada keempat persona. Akurasi persona Islam GPT-OSS meningkat
2,67 poin, sedangkan akurasi persona Kristen Llama 3.1 menurun 4,47 poin.
Pengulangan Yom Kippur pada satu respons Qwen persona Yahudi tidak disertai
perbedaan akurasi I2--E yang signifikan pada model dan persona tersebut.

I3 menyampaikan persona melalui masjid, gereja, sinagoge atau Shabbat, serta
tidak adanya komunitas agama pada persona ateis. Pada Ministral 8B, I3 lebih
sering membagi respons menjadi langkah dan daftar berpoin serta memiliki
variasi kata (TTR-300) yang lebih rendah daripada E. Perubahan tersebut
disertai penurunan akurasi 2,88--7,05 poin pada keempat persona. Pada skenario
yang sama, akurasi persona Islam GPT-OSS meningkat 2,37 poin, sedangkan
akurasi persona Kristen Qwen menurun 1,21 poin.

DeepSeek R1-Distill 32B dan Llama 3.3 70B mengalami perubahan bentuk respons
pada I1, I2, dan I3 tanpa perbedaan akurasi I--E yang signifikan. DeepSeek
menulis respons implisit yang lebih singkat dan mengubah penempatan persamaan,
sedangkan Llama 3.3 lebih sering membagi respons menjadi langkah dan memakai
kata penghubung sebab-akibat. Perubahan bahasa pada kedua model tersebut tidak
diikuti perubahan akurasi setelah koreksi Holm.

Perbandingan antarpersona juga lebih sering berbeda pada bahasa
respons daripada akurasi. Penanda linguistik umum berbeda pada seluruh model dan bahasa
keagamaan berbeda pada E, tetapi hanya empat dari 144 perbandingan akurasi
antarpersona yang signifikan. Keempatnya melibatkan persona Islam dengan
akurasi lebih rendah pada Ministral 8B atau Llama 3.1 8B. Persona Islam tidak
selalu memperoleh akurasi paling rendah pada model lain.

Persona agama mengubah kosakata dan bentuk respons pada keenam model, sedangkan
perubahan akurasi hanya ditemukan pada empat model. E paling kuat membedakan
bahasa keagamaan karena identitas disebutkan secara langsung. I1, I2, dan I3
lebih banyak membedakan susunan respons, tetapi kenaikan atau penurunan akurasi
tetap mengikuti model dan persona yang digunakan. Pada GSM8K, pengaruh persona
agama karena itu lebih konsisten pada bahasa respons daripada akurasi.

\subsubsection{Ringkasan Domain Agama}
\label{subsubsec:ringkasan-domain-agama-gsm8k}

Pada domain agama, bagian respons yang berubah berbeda menurut cara persona
disampaikan. Skenario E lebih banyak memunculkan istilah yang sesuai dengan
identitas agama karena identitas pengguna disebut secara langsung. Pada I1,
I2, dan I3, istilah agama hampir tidak muncul, tetapi susunan jawaban tetap
berubah setelah percakapan \textit{warmup}. Perubahan pada E lebih banyak
terdapat pada kosakata identitas, sedangkan perubahan pada skenario implisit
lebih banyak terdapat pada bentuk umum jawaban.

Perubahan akurasi berbeda antarmodel. Ministral 8B mengalami penurunan
2,58--8,79 poin pada ketiga skenario implisit untuk seluruh persona agama.
GPT-OSS 20B mengalami kenaikan 2,06--2,67 poin pada persona Islam, sedangkan
Llama 3.1 8B dan Qwen3 Next 80B-A3B mengalami penurunan pada sebagian
skenario persona Kristen. DeepSeek R1-Distill 32B dan Llama 3.3 70B tidak
menghasilkan perbedaan akurasi I--E yang signifikan.

Perbandingan antarpersona hanya menghasilkan empat perbedaan akurasi dari 144
perbandingan. Keempatnya menunjukkan akurasi persona Islam yang lebih rendah:
tiga pada Ministral 8B terhadap persona ateis atau Yahudi dan satu pada Llama
3.1 8B terhadap persona Kristen. Kecenderungan ini hanya berlaku pada empat
hasil signifikan tersebut; tidak ada satu persona agama yang secara konsisten
memperoleh akurasi paling tinggi atau paling rendah pada keenam model.
Identitas agama lebih konsisten membedakan kosakata dan cara model menulis
jawaban daripada akurasi GSM8K.

\subsection{Domain Regional-Kultural}
\label{subsec:hasil-regional-gsm8k}

Domain regional-kultural mencakup persona Amerika Utara, Asia Timur, Asia
Selatan, dan Afrika Sub-Sahara. Setiap persona diuji dengan satu skenario
eksplisit dan tiga skenario implisit. I1 menggunakan nama dan situasi
pengambilan keputusan, I2 menggunakan hari raya atau perayaan budaya, dan I3
menggunakan tradisi pernikahan. Ketiganya dianalisis secara terpisah karena
menggunakan \textit{warmup} yang berbeda. Analisis ini menggunakan 134.538
respons dari 1.319 soal GSM8K.

\subsubsection{Akurasi Respons}
\label{subsubsec:akurasi-regional-gsm8k}

\par\noindent\textbf{a. Akurasi Setiap Skenario}\par

Tabel~\ref{tab:akurasi-regional-gsm8k} menunjukkan akurasi keenam model pada
skenario netral serta skenario yang menggunakan persona Amerika Utara, Asia
Timur, Asia Selatan, dan Afrika Sub-Sahara.

\begin{table}[H]
    \centering
    \caption{Akurasi GSM8K pada domain regional-kultural menurut persona dan skenario (\%)}
    \label{tab:akurasi-regional-gsm8k}
    \scriptsize
    \setlength{\tabcolsep}{3pt}
    \renewcommand{\arraystretch}{1.08}
    \begin{tabularx}{\textwidth}{
        >{\raggedright\arraybackslash}p{1.75cm}
        >{\centering\arraybackslash}p{0.65cm}
        *{6}{>{\centering\arraybackslash}X}}
        \toprule
        \textbf{Persona} & \textbf{Skenario} & \textbf{Qwen3} &
        \textbf{Llama 3.1} & \textbf{Ministral} & \textbf{DeepSeek} &
        \textbf{GPT-OSS} & \textbf{Llama 3.3} \\
        \midrule
        Netral & N & 94,08 & 84,43 & 93,47 & 93,93 & 93,93 & 91,96 \\
        \addlinespace[2pt]
        Amerika Utara & E  & 96,59 & 82,64 & 91,96 & 94,62 & 92,02 & 92,27 \\
                       & I1 & 95,07 & 82,87 & 91,96 & 93,10 & 91,73 & 92,04 \\
                       & I2 & 95,22 & 79,98 & 92,34 & 93,48 & 92,33 & 92,95 \\
                       & I3 & 95,83 & 82,94 & 91,28 & 92,57 & 90,52 & 92,87 \\
        \addlinespace[2pt]
        Asia Timur & E  & 95,91 & 83,24 & 91,05 & 94,62 & 92,94 & 93,25 \\
                   & I1 & 95,30 & 81,20 & 91,43 & 93,78 & 92,49 & 93,33 \\
                   & I2 & 95,68 & 82,26 & 90,30 & 93,86 & 92,49 & 91,74 \\
                   & I3 & 95,75 & 81,35 & 91,58 & 94,77 & 92,41 & 92,27 \\
        \addlinespace[2pt]
        Asia Selatan & E  & 96,44 & 83,70 & 89,69 & 95,00 & 91,79 & 94,16 \\
                     & I1 & 95,30 & 82,94 & 90,30 & 93,56 & 92,26 & 93,33 \\
                     & I2 & 95,07 & 81,80 & 91,58 & 93,86 & 92,26 & 92,95 \\
                     & I3 & 95,00 & 83,47 & 89,99 & 94,47 & 91,12 & 92,34 \\
        \addlinespace[2pt]
        Afrika Sub-Sahara & E  & 96,51 & 82,64 & 91,36 & 94,54 & 92,12 & 92,95 \\
                          & I1 & 95,75 & 80,97 & 90,83 & 93,40 & 91,66 & 92,19 \\
                          & I2 & 94,77 & 80,59 & 88,17 & 93,03 & 91,65 & 93,40 \\
                          & I3 & 95,52 & 80,44 & 91,13 & 94,16 & 92,34 & 92,19 \\
        \bottomrule
    \end{tabularx}
\end{table}

Pada domain regional-kultural, perubahan akurasi berbeda antarmodel dan paling
besar terdapat pada Ministral 8B. Selisih antara akurasi tertinggi dan terendah
dalam satu model berkisar 2,42--5,30 poin persentase. Selisih paling kecil
terdapat pada Llama 3.3 70B.
Gambar~\ref{fig:akurasi-absolut-regional-gsm8k} menunjukkan perbandingan hasil
pada skenario netral, skenario eksplisit, dan rata-rata I1, I2, dan I3.

\begin{figure}[H]
    \centering
    \includegraphics[
        width=\textwidth,
        height=0.46\textheight,
        keepaspectratio
    ]{image/bab6/gsm8k/regional_kultural/akurasi/12_regional_absolute_accuracy_explicit_implicit.png}
    \caption[Akurasi GSM8K menurut skenario pada domain regional-kultural]
    {Perbandingan akurasi GSM8K antara skenario netral, skenario eksplisit,
    dan rata-rata skenario implisit pada domain regional-kultural}
    \label{fig:akurasi-absolut-regional-gsm8k}
\end{figure}

Seluruh skenario persona pada Qwen3 Next 80B-A3B memperoleh akurasi lebih
tinggi daripada hasil pada skenario netral, sedangkan seluruh skenario persona
pada Llama 3.1 8B, Ministral 8B, dan GPT-OSS 20B memperoleh akurasi lebih
rendah. Pada Llama 3.3 70B, hanya I2 persona Asia Timur yang memperoleh akurasi
lebih rendah. Pada DeepSeek R1-Distill 32B, seluruh E memperoleh akurasi lebih
tinggi daripada hasil pada skenario netral. Pada model tersebut, I1 dan I2 pada
keempat persona lebih rendah, sedangkan I3 lebih rendah hanya pada persona
Amerika Utara. Penurunan paling besar terdapat pada I2 persona Afrika
Sub-Sahara pada Ministral 8B, yaitu 5,30 poin persentase.

Gambar~\ref{fig:delta-implisit-eksplisit-regional-gsm8k} menunjukkan selisih
rata-rata I1, I2, dan I3 terhadap E untuk setiap persona. Nilai positif
menunjukkan rata-rata akurasi skenario implisit yang lebih tinggi daripada E.
Nilai negatif menunjukkan rata-rata akurasi skenario implisit yang lebih
rendah daripada E.

\begin{figure}[H]
    \centering
    \includegraphics[
        width=\textwidth,
        height=0.46\textheight,
        keepaspectratio
    ]{image/bab6/gsm8k/regional_kultural/akurasi/13_regional_delta_implicit_vs_explicit.png}
    \caption[Selisih akurasi antarskenario pada domain regional-kultural]
    {Selisih rata-rata akurasi antara skenario implisit dan eksplisit pada
    domain regional-kultural}
    \label{fig:delta-implisit-eksplisit-regional-gsm8k}
\end{figure}

Pada Qwen3 Next 80B-A3B, rata-rata akurasi skenario implisit lebih rendah
0,33--1,31 poin daripada E pada keempat persona. Penurunan pada Llama 3.1 8B
berkisar 0,71--1,97 poin, sedangkan pada DeepSeek R1-Distill 32B berkisar
0,48--1,57 poin. Selisih pada GPT-OSS 20B berada pada rentang $-0{,}49$ sampai
$+0{,}08$ poin. Ministral 8B memperoleh selisih $-1{,}31$ sampai $+0{,}94$
poin, sedangkan Llama 3.3 70B memperoleh selisih $-1{,}29$ sampai $+0{,}35$
poin.

\par\noindent\textbf{b. Hasil Perbandingan Statistik}\par
\label{subsubsec:perbandingan-regional-gsm8k}

Uji Q Cochran digunakan untuk membandingkan akurasi E, I1, I2, dan I3 pada
model dan persona regional-kultural yang sama. Exact McNemar digunakan untuk
membandingkan I1, I2, dan I3 dengan E, membandingkan E dan ketiga skenario
implisit dengan N, serta membandingkan persona regional-kultural pada model dan
skenario yang sama. Hasil I--E dimasukkan ke pembahasan utama jika Exact
McNemar setelah koreksi Holm dan uji Q Cochran pada kelompok yang sama
signifikan. Pengujian mencakup 72 perbandingan I--E, 96 perbandingan dengan N,
dan 144 perbandingan antarpersona regional-kultural.

Uji Q Cochran signifikan pada enam dari 24 kelompok model dan persona. Tiga
kelompok terdapat pada persona Amerika Utara, yaitu Qwen3 Next 80B-A3B, Llama
3.1 8B, dan DeepSeek R1-Distill 32B. Tiga kelompok lainnya terdapat pada Qwen3
Next 80B-A3B persona Asia Selatan dan Afrika Sub-Sahara serta Ministral 8B
persona Afrika Sub-Sahara.

Setelah koreksi Holm, 10 perbandingan I--E, 15
perbandingan dengan N, dan tiga perbandingan antarpersona tetap signifikan.
Satu hasil I--E, yaitu I3 persona Asia Selatan pada Llama 3.3 70B, memiliki
$p_{\mathrm{Holm}}=0{,}045$, tetapi kelompoknya tidak didukung Q Cochran dengan
$p=0{,}077$. Oleh karena itu, sembilan perbandingan I--E dimasukkan ke hasil
utama.

Dari 96 perbandingan dengan skenario netral, 15 tetap signifikan setelah
koreksi Holm. Enam hasil terdapat pada Ministral 8B, empat pada Llama 3.1 8B,
tiga pada Qwen3 Next 80B-A3B, dan dua pada GPT-OSS 20B. Ketiga hasil pada Qwen3
berupa kenaikan akurasi E, sedangkan 12 hasil pada tiga model lainnya berupa
penurunan.

Kesembilan perbandingan I--E yang didukung Q Cochran dan koreksi Holm berupa
penurunan akurasi. Tujuh hasil terdapat pada Qwen3 Next 80B-A3B, yaitu I1 dan
I2 Amerika Utara, I1--I3 Asia Selatan, serta I2 dan I3 Afrika Sub-Sahara,
dengan penurunan 1,06--1,74 poin. Pada DeepSeek R1-Distill 32B, I3 Amerika
Utara lebih rendah 2,05 poin daripada E. Penurunan terbesar terdapat pada I2
Afrika Sub-Sahara pada Ministral 8B sebesar 3,18 poin. Persona Asia Timur tidak
memiliki hasil I--E yang didukung kedua uji.

Tiga dari 144 perbandingan antarpersona regional-kultural tetap signifikan
setelah koreksi Holm. Pada Ministral 8B, akurasi persona Amerika Utara pada I2
lebih tinggi 4,17 poin daripada Afrika Sub-Sahara, sedangkan persona Asia
Selatan lebih tinggi 3,41 poin daripada Afrika Sub-Sahara pada skenario yang
sama. Pada DeepSeek R1-Distill 32B, akurasi persona Amerika Utara pada I3 lebih
rendah 2,20 poin daripada Asia Timur. Ketiga hasil tersebut hanya terdapat pada
I2 dan I3 pada Ministral 8B dan DeepSeek R1-Distill 32B.

Gambar~\ref{fig:effect-regional-gsm8k} menunjukkan \textit{effect size}
Cohen's $h$ pada 52 jenis perbandingan per model. Panel pertama dan kedua
masing-masing memuat 14 perbandingan antarskenario, sedangkan panel ketiga dan
keempat masing-masing memuat 12 perbandingan antarpersona. Warna biru
menunjukkan akurasi skenario atau persona pertama lebih rendah daripada
pembanding, sedangkan warna merah menunjukkan akurasi yang lebih tinggi. Tanda
bintang menunjukkan hasil yang tetap signifikan setelah koreksi Holm, tanda
$\dagger$ menunjukkan perbandingan dengan N, dan tanda $\ddagger$ menunjukkan
perbandingan antarpersona. Tanda bintang pada I3 persona Asia Selatan pada
Llama 3.3 70B hanya didukung koreksi Holm sehingga tidak dimuat dalam tabel;
nilai \textit{p-value} lengkap terdapat pada
Lampiran~\ref{chap:lampiran-statistik-regional}.

\begin{figure}[H]
    \centering
    \includegraphics[
        width=\linewidth,
        height=0.84\textheight,
        keepaspectratio,
        trim=0 90bp 0 0,
        clip
    ]{image/bab6/gsm8k/regional_kultural/akurasi/27_regional_all_52_effect_holm_stacked_26_26.png}
    \caption[Ukuran efek dan signifikansi akurasi domain regional-kultural]
    {Ukuran efek Cohen's $h$ dan hasil uji Exact McNemar setelah koreksi Holm
    pada 52 jenis perbandingan per model dalam domain regional-kultural.\\[-2pt]
    {\footnotesize\textit{Catatan:} * = signifikan setelah koreksi Holm
    ($p_{\mathrm{Holm}} < 0{,}05$); $\dagger$ = perbandingan terhadap
    skenario netral bersama (analisis sensitivitas; protokol berbeda);
    $\ddagger$ = perbandingan antarkawasan pada E, I1, I2, dan I3.
    E = eksplisit; I1--I3 = implisit;
    N = netral. NA = Amerika Utara; EA = Asia Timur; SA = Asia Selatan;
    SSA = Afrika Sub-Sahara. Cohen's $h$ dihitung dari skenario pertama
    dikurangi skenario pembanding.}}
    \label{fig:effect-regional-gsm8k}
\end{figure}

Tabel~\ref{tab:hasil-signifikan-regional-gsm8k} memuat sembilan perbandingan
I--E yang didukung Q Cochran dan koreksi Holm. Hasil I1, I2, dan I3 pada model
dan persona yang sama digabungkan dalam satu baris jika seluruh akurasinya
menurun. Kolom selisih dan Cohen's $h$ menunjukkan rentang hasil dari skenario
yang tercantum.

\begin{table}[H]
    \centering
    \caption[Perbandingan akurasi I--E pada domain regional-kultural]
    {Perbandingan akurasi I--E pada domain regional-kultural yang didukung oleh
    uji Q Cochran dan tetap signifikan setelah koreksi Holm}
    \label{tab:hasil-signifikan-regional-gsm8k}
    \footnotesize
    \renewcommand{\arraystretch}{1.15}
    \begin{adjustbox}{max width=\textwidth}
    \begin{tabular}{
        >{\raggedright\arraybackslash}p{2.4cm}
        >{\raggedright\arraybackslash}p{3.4cm}
        >{\centering\arraybackslash}p{1.3cm}
        >{\centering\arraybackslash}p{2.6cm}
        >{\centering\arraybackslash}p{2.2cm}}
        \toprule
        \textbf{Persona} & \textbf{Model} & \textbf{Skenario} &
        \textbf{Selisih I--E} & \textbf{Rentang $h$} \\
        \midrule
        Amerika Utara & Qwen3 Next 80B-A3B & I1--I2 &
        $-1{,}52$ s.d. $-1{,}36$ & $-0{,}08$ s.d. $-0{,}07$ \\
        Amerika Utara & DeepSeek R1-Distill 32B & I3 &
        $-2{,}05$ & $-0{,}08$ \\
        Asia Selatan & Qwen3 Next 80B-A3B & I1--I3 &
        $-1{,}44$ s.d. $-1{,}14$ & $-0{,}07$ s.d. $-0{,}06$ \\
        Afrika Sub-Sahara & Ministral 8B & I2 &
        $-3{,}18$ & $-0{,}11$ \\
        Afrika Sub-Sahara & Qwen3 Next 80B-A3B & I2--I3 &
        $-1{,}74$ s.d. $-1{,}06$ & $-0{,}09$ s.d. $-0{,}05$ \\
        \bottomrule
    \end{tabular}
    \end{adjustbox}
\end{table}

Nilai $|h|$ pada sembilan perbandingan I--E berkisar 0,05--0,11 dan seluruhnya
berada di bawah 0,20. Penurunan terbesar terdapat pada I2 persona Afrika
Sub-Sahara pada Ministral 8B. Tujuh penurunan terdapat pada Qwen3 Next
80B-A3B, sedangkan dua lainnya terdapat pada DeepSeek R1-Distill 32B dan
Ministral 8B.


\par\noindent\textbf{c. Penggunaan Token}\par
\label{subsubsec:token-regional-gsm8k}

Ketiga skenario implisit memberi konteks regional-kultural yang berbeda: I1
menggunakan nama dan situasi pengambilan keputusan, I2 menggunakan hari raya
atau perayaan budaya, dan I3 menggunakan tradisi pernikahan. Seluruh percakapan
\textit{warmup} dikirim bersama soal sehingga rata-rata \textit{input token}
I1--I3 meningkat 278,2--1.791,7\% dibandingkan dengan E. Skenario dengan
\textit{input token} tertinggi berbeda menurut model dan persona karena panjang
pertanyaan serta respons \textit{warmup} juga berbeda.

Penggunaan \textit{completion token} berbeda antarmodel. Pada DeepSeek
R1-Distill 32B, rata-rata I1--I3 meningkat 69,1--74,8\% dibandingkan dengan E.
Qwen3 Next 80B-A3B mengalami kenaikan 6,3--11,7\%, sedangkan GPT-OSS 20B
mengalami penurunan 12,0--15,5\%. Pada Llama 3.1 8B, perubahannya berkisar
$-8{,}4$ sampai $+9{,}8$\%, pada Ministral 8B meningkat 2,9--10,1\%, dan pada
Llama 3.3 70B berkisar $-0{,}6$ sampai $+5{,}9$\%. Rinciannya terdapat pada
Gambar~\ref{fig:lampiran-token-regional-implisit-eksplisit} di
Lampiran~\ref{chap:lampiran-statistik-regional}.

Jika \textit{input token} dan \textit{completion token} dijumlahkan, rata-rata
total token pada I1--I3 lebih tinggi 79,1--330,1\% daripada E. Kenaikan
tertinggi terdapat pada persona Asia Timur pada Ministral 8B, sedangkan kenaikan
terendah terdapat pada persona Asia Selatan pada Qwen3 Next 80B-A3B. Jumlah
tertinggi terdapat pada persona Amerika Utara I2 pada Qwen3 Next 80B-A3B, yaitu
rata-rata sekitar 4.721 token per soal. Pada persona Asia Selatan, persentase
kenaikan Qwen lebih rendah karena E sudah menggunakan sekitar 1.071 total token
per soal.

Perubahan \textit{completion token} dan akurasi berbeda antarmodel. Pada
DeepSeek R1-Distill 32B, kenaikan \textit{completion token} sebesar
69,1--74,8\% disertai penurunan akurasi sebesar 0,48--1,57 poin. Pada Qwen3
Next 80B-A3B, kenaikan 6,3--11,7\% juga disertai penurunan akurasi
0,33--1,31 poin. GPT-OSS 20B menggunakan 12,0--15,5\% lebih sedikit
\textit{completion token}, sedangkan perubahan akurasinya hanya berkisar
$-0{,}49$ sampai $+0{,}08$ poin. Pada Ministral 8B, kenaikan
\textit{completion token} sebesar 2,9--10,1\% disertai perubahan akurasi antara
$-1{,}31$ dan $+0{,}94$ poin.

Pada rata-rata I1--I3, selisih total token antara persona dengan jumlah
tertinggi dan terendah dalam satu model berkisar 10,2--66,5\%. Persona Amerika
Utara menggunakan total token tertinggi pada Llama 3.1 8B, Llama 3.3 70B, dan
Qwen3 Next 80B-A3B. Jumlah tertinggi terdapat pada persona Asia Timur untuk
Ministral 8B dan GPT-OSS 20B serta persona Afrika Sub-Sahara untuk DeepSeek
R1-Distill 32B. Selisih persentase terbesar terdapat pada Ministral 8B, yaitu
66,5\% antara persona Asia Timur dan Asia Selatan. Pada Qwen3 Next 80B-A3B,
selisih kedua persona tersebut mencapai 987 token dan terutama berasal dari
\textit{input token}.

Perbandingan penggunaan token dan akurasi juga berbeda menurut persona. Pada Qwen3
Next 80B-A3B, persona Amerika Utara menggunakan
rata-rata total token tertinggi, sedangkan akurasi I1 dan I2 lebih rendah
1,36--1,52 poin daripada E. Pada Ministral 8B, persona Asia Timur menggunakan
rata-rata total token tertinggi, sedangkan perbandingan antarkawasan pada I2
menempatkan persona Amerika Utara dan Asia Selatan masing-masing 4,17 dan 3,41
poin lebih tinggi daripada persona Afrika Sub-Sahara. Pada DeepSeek R1-Distill
32B, persona Afrika Sub-Sahara menggunakan
rata-rata total token tertinggi, sedangkan perbandingan antarkawasan yang signifikan
terdapat pada I3 dengan akurasi Amerika Utara 2,20 poin lebih rendah daripada
Asia Timur. Pada ketiga model tersebut, rata-rata total token tertinggi tidak
disertai keunggulan akurasi yang signifikan pada persona yang sama.

\subsubsection{Penanda Linguistik}
\label{subsubsec:penanda-regional-gsm8k}

Pada domain regional-kultural, penanda linguistik dibandingkan melalui tiga
jenis perbandingan. Skenario E, I1, I2, dan I3 dibandingkan dengan N untuk
melihat perubahan dari skenario tanpa persona. Skenario I1, I2, dan I3 juga
dibandingkan dengan E pada persona yang sama untuk melihat perbedaan respons
setelah skenario implisit digunakan. Keempat persona regional-kultural
kemudian dibandingkan dalam enam pasangan pada soal dan skenario yang sama.

Analisis memisahkan 16 penanda linguistik umum, kategori pembuka, dua penanda
linguistik bahasa matematis, 14 penanda linguistik regional-kultural, istilah
hasil audit, dan akurasi.
Pemisahan ini diperlukan karena penyebutan kawasan tidak sama dengan perubahan
cara model menulis penyelesaian matematika. Respons yang dianalisis hanya teks
yang terlihat oleh pengguna; penalaran tersembunyi tidak ikut dihitung.

I1 membawa percakapan tentang keputusan pribadi atau sosial, I2 membawa
perayaan dan tradisi hari besar, sedangkan I3 membawa adat pernikahan. Ketiga
skenario tersebut tetap dihitung terpisah karena isi \textit{warmup}-nya tidak
setara. Respons \textit{warmup} menjadi riwayat percakapan, tetapi tidak masuk
ke dalam hitungan penanda linguistik atau akurasi jawaban GSM8K.

Penanda linguistik numerik diuji dengan Friedman untuk perbandingan kelompok dan
Wilcoxon \textit{signed-rank} untuk pasangan, lalu nilai-p pasangan dikoreksi
dengan Holm. Untuk setiap model, persona, dan penanda linguistik, keluarga Holm memuat
empat perbandingan terhadap N serta tiga perbandingan I1--I3 terhadap E.
Penanda linguistik khusus regional-kultural yang berupa ada atau tidak ada
diuji dengan Exact McNemar, sedangkan komposisi pembuka diuji dengan
Stuart--Maxwell. Akurasi memakai Q Cochran sebagai uji kelompok dan Exact
McNemar untuk pasangan seperti pada bagian evaluasi akurasi.

\par\noindent\textbf{a. Perubahan terhadap Skenario Netral}\par

Gambar~\ref{fig:marker-regional-terhadap-netral} membandingkan E, I1, I2, dan
I3 dengan N pada persona regional-kultural yang sama. Empat baris memisahkan
hasil Amerika Utara, Asia Timur, Asia Selatan, dan Afrika Sub-Sahara, sedangkan
empat kolom menunjukkan E--N, I1--N, I2--N, dan I3--N. Satu titik mewakili
hasil satu penanda linguistik pada satu model, warna titik membedakan kelompok
penanda linguistik pada Tabel~\ref{tab:aturan-penanda}, dan ukuran titik
mengikuti jumlah pasangan respons yang memiliki nilai berbeda. Sumbu horizontal
menunjukkan \textit{effect size} \textit{rank-biserial}; nilai positif
menunjukkan penanda linguistik lebih tinggi pada skenario persona, sedangkan
nilai negatif menunjukkan penanda linguistik lebih tinggi pada N. Sumbu
vertikal menunjukkan $-\log_{10}(p_{\text{Holm}})$, garis putus-putus menandai
batas signifikansi setelah koreksi Holm, dan nilai di atas 50 ditempatkan pada
50.

\begin{landscape}
\begin{figure}[H]
    \centering
    \includegraphics[
        width=\linewidth,
        height=0.82\textheight,
        keepaspectratio
    ]{image/BabVI/gsm8k/regional_kultural/14_numeric_effect_volcano_vs_neutral.png}
    \caption{Ukuran efek penanda linguistik terhadap skenario netral pada domain regional-kultural}
    \label{fig:marker-regional-terhadap-netral}
\end{figure}
\end{landscape}

Ringkasan penanda linguistik dengan ukuran efek terbesar terdapat pada
Lampiran~\ref{tab:lampiran-penanda-regional-ringkas}.
Pada keenam model, E, I1, I2, dan I3 menghasilkan kalimat yang lebih panjang
serta lebih sedikit menggunakan daftar berpoin dan subjudul daripada N
pada keempat persona regional-kultural. Perubahan panjang kalimat, daftar
berpoin, dan subjudul tetap signifikan setelah koreksi Holm pada seluruh
model, persona, dan skenario. Langkah eksplisit juga lebih sedikit pada seluruh
perbandingan; hanya I2 persona Afrika Sub-Sahara pada GPT-OSS 20B yang tidak
tetap signifikan setelah koreksi Holm.

Pada DeepSeek R1-Distill 32B, ketiga skenario implisit menghasilkan lebih
sedikit kata dan LaTeX \textit{display} daripada N pada keempat persona. Kata
penghubung sebab-akibat lebih banyak digunakan, dan ketiga perubahan tersebut
tetap signifikan setelah koreksi Holm. Pada Qwen3 Next 80B-A3B, E dan ketiga
skenario implisit menghasilkan lebih sedikit kata serta lebih sedikit
menggunakan kata penghubung sebab-akibat
daripada N pada keempat persona. LaTeX \textit{display} juga lebih sedikit
digunakan pada ketiga skenario implisit, dan seluruh perubahan tersebut tetap
signifikan setelah koreksi Holm.

Pada Llama 3.3 70B, jumlah kata tidak berbeda secara konsisten dari N pada
keempat persona. Kata penghubung sebab-akibat lebih banyak digunakan pada I1
dan I2 serta tetap signifikan setelah koreksi Holm untuk keempat persona.
Kalimat juga lebih panjang dan penanda linguistik struktur berkurang pada keempat skenario
persona. Pada GPT-OSS 20B, LaTeX \textit{display} lebih sedikit digunakan pada
E, I1, I2, dan I3 untuk keempat persona serta tetap signifikan setelah koreksi
Holm. Jumlah kata pada skenario implisit tidak menghasilkan pola signifikan
yang sama pada keempat persona.

Pada Llama 3.1 8B, I1 menghasilkan lebih banyak kata daripada N pada keempat
persona dan tetap signifikan setelah koreksi Holm, sedangkan I2 dan I3 tidak
menghasilkan pola signifikan yang sama. Kalimat tetap lebih panjang serta langkah
eksplisit, daftar berpoin, dan subjudul lebih sedikit digunakan daripada N.
Pada Ministral 8B, ketiga skenario implisit menghasilkan lebih banyak kata dan
kalimat yang lebih panjang, disertai penggunaan penanda linguistik struktur yang lebih
sedikit. Kenaikan jumlah kata pada I1 signifikan untuk keempat persona, pada I2
untuk tiga persona, dan pada I3 untuk persona Amerika Utara serta Asia Timur.

Kalimat yang lebih panjang muncul bersama berkurangnya langkah eksplisit,
daftar berpoin, dan subjudul. DeepSeek R1-Distill 32B dan Qwen3 Next
80B-A3B menggunakan lebih sedikit kata, sedangkan Ministral 8B menggunakan
lebih banyak kata. Kenaikan rata-rata panjang kalimat karena itu tidak selalu
disertai kenaikan jumlah kata.

Pola susunan ini muncul pada keempat kawasan dan
lebih banyak berbeda menurut model serta skenario daripada menurut satu kawasan. E
menambahkan persona secara langsung, sedangkan I1, I2, dan I3 menambahkan
percakapan \textit{warmup} sebelum soal. Perbandingan I--N karena itu mencakup
cara penyampaian persona dan penambahan riwayat percakapan karena N tidak
menggunakan \textit{warmup} yang setara.

\par\noindent\textbf{b. Perbandingan Antarkawasan}\par

Perbandingan antarkawasan dilakukan pada model, skenario, dan soal yang sama.
Enam pasangan yang dibandingkan adalah Amerika Utara--Asia Timur, Amerika
Utara--Asia Selatan, Amerika Utara--Afrika Sub-Sahara, Asia Timur--Asia
Selatan, Asia Timur--Afrika Sub-Sahara, dan Asia Selatan--Afrika Sub-Sahara.
Hasil setiap pasangan tetap dilaporkan terpisah dan tidak digabung menjadi
satu kelompok kawasan.

Setiap pasangan diuji pada 14 penanda linguistik khusus regional-kultural. Keluarga
koreksi Holm memuat 24 perbandingan antarkawasan per model dan penanda linguistik, yaitu enam
pasangan kawasan pada E, I1, I2, dan I3. Setelah koreksi Holm pada uji Exact
McNemar, E memiliki 43 kombinasi model, pasangan kawasan, dan penanda linguistik dengan
$p_{\text{Holm}}<0{,}05$. I2 memiliki dua, I3 memiliki lima, sedangkan seluruh
perbandingan I1 memperoleh $p_{\text{Holm}}\geq 0{,}05$.

Gambar~\ref{fig:marker-regional-antarkawasan-e} memakai selisih prevalensi
dalam poin persentase, yaitu kawasan pertama dikurangi kawasan kedua. Nilai
positif menandai penanda linguistik yang lebih sering muncul pada kawasan pertama. Titik
berwarna mewakili model, garis horizontal merupakan
interval kepercayaan 95\%, dan berlian hitam merupakan median keenam model.
Titik yang lebih pucat tidak signifikan setelah koreksi Holm.

\begin{figure}[H]
    \centering
    \includegraphics[
        width=\textwidth,
        height=0.78\textheight,
        keepaspectratio
    ]{image/BabVI/gsm8k/regional_kultural/09_direct_identity_regional_forest_explicit.png}
    \caption[Selisih penanda linguistik antarkawasan]
    {Selisih prevalensi penanda linguistik pada perbandingan antarkawasan
    dalam skenario E}
    \label{fig:marker-regional-antarkawasan-e}
\end{figure}

\begin{table}[H]
    \centering
    \caption{Hasil signifikan penanda linguistik khusus regional-kultural antarkawasan}
    \label{tab:penanda-khusus-regional-antarkawasan}
    \footnotesize
    \renewcommand{\arraystretch}{1.12}
    \begin{adjustbox}{max width=\textwidth}
    \begin{tabular}{lcll}
        \toprule
        \textbf{Model} & \textbf{Jumlah} & \textbf{Skenario} & \textbf{Selisih absolut terbesar} \\
        \midrule
        DeepSeek R1-Distill 32B & 0 & -- & -- \\
        Llama 3.1 8B & 21 & E & Label Asia Timur: 50,57 poin \\
        Llama 3.3 70B & 29 & E, I2, I3 & Sapaan formulaik: 14,88 poin \\
        Ministral 8B & 0 & -- & -- \\
        GPT-OSS 20B & 0 & -- & -- \\
        Qwen3 Next 80B-A3B & 0 & -- & -- \\
        \bottomrule
    \end{tabular}
    \end{adjustbox}
\end{table}

Tabel~\ref{tab:penanda-khusus-regional-antarkawasan} menghitung hasil
signifikan untuk setiap kombinasi pasangan kawasan, skenario, dan penanda linguistik.
Seluruh 43 hasil E hanya terdapat pada Llama 3.1 8B dan Llama 3.3 70B,
masing-masing sebanyak 21 dan 22. Pada Llama 3.1 8B, selisih penyebutan label
kawasan mencapai 50,57 poin untuk Asia Timur, 44,35 poin untuk Asia Selatan,
dan 30,33 poin untuk Afrika Sub-Sahara. Pada Llama 3.3 70B, selisih terbesar
sebesar 14,88 poin terdapat pada sapaan formulaik.

Contohnya terdapat pada Llama 3.1 8B untuk soal \texttt{gsm8k\_00265}.
Respons E persona Amerika Utara dimulai dengan ``\textit{Let's break it down
step by step}'', sedangkan respons E persona Asia Selatan dimulai dengan
``\textit{Namaste! Let's solve the problem step by step.}'' Kedua respons memperoleh
angka GSM8K 400 dan dinilai benar. Contoh ini memperlihatkan penanda linguistik sapaan
yang dihitung tanpa mengaitkannya dengan perbedaan akurasi.

Pada Llama 3.3 70B, dua hasil I2 terdapat pada tradisi atau objek bernama dengan
selisih 1,06 poin. Lima hasil I3 pada model yang sama terdapat pada rujukan
budaya umum dan tradisi atau objek bernama dengan selisih 0,91--0,99 poin.
Ketujuh hasil skenario implisit ditemukan pada satu model dengan selisih tidak
melebihi 1,06 poin. Perbedaan antarkawasan pada skenario implisit lebih jarang
dan lebih kecil daripada perbedaan E. Hasil lengkap E, I1, I2, dan I3 terdapat pada
Gambar~\ref{fig:lampiran-marker-langsung-regional-e} sampai
Gambar~\ref{fig:lampiran-marker-langsung-regional-i3} di Lampiran.

\par\noindent\textbf{c. Bentuk Respons dan Bahasa Matematis}\par

Perubahan bentuk respons perlu dipisahkan dari penggunaan istilah budaya.
Jumlah kata, rata-rata panjang kalimat, variasi kata, langkah eksplisit, dan
bahasa matematis mengukur bentuk penyelesaian yang ditulis model. Penyebutan
kawasan, sapaan, dan istilah budaya mengukur penanda linguistik identitas pengguna yang
muncul pada respons.

Pada DeepSeek R1-Distill 32B, I1, I2, dan I3 menghasilkan lebih sedikit kata,
kalimat yang lebih pendek, dan lebih sedikit LaTeX \textit{display} daripada E
pada keempat persona. Variasi kata, kata penghubung sebab-akibat, dan operasi
aritmetika di dalam kalimat justru meningkat, dan seluruh perubahan tersebut
tetap signifikan setelah koreksi Holm. Pada Ministral 8B, ketiga skenario
implisit justru menghasilkan lebih banyak kata dan kalimat yang lebih panjang
daripada E pada keempat persona, sedangkan variasi kata menurun; ketiga
perubahan tersebut juga tetap signifikan setelah koreksi Holm. Selisih akurasi
signifikan terdapat pada I3 persona
Amerika Utara di DeepSeek sebesar 2,05 poin dan I2 persona Afrika Sub-Sahara
di Ministral sebesar 3,18 poin. Bentuk perubahan bahasa dan akurasi berbeda
menurut model.

Perbedaan bentuk respons terlihat pada DeepSeek R1-Distill 32B untuk
\texttt{gsm8k\_01163} persona Amerika Utara. Respons E membuka dengan
``\textit{Solution: Let's break down the problem step by step}'', memakai
subjudul untuk dua tahap perhitungan, lalu menulis persamaan dalam LaTeX
\textit{display}. Respons I3 hanya menulis ``\textit{Geb is 3 years old.}''
Keduanya dinilai benar. Pada keseluruhan soal DeepSeek untuk persona dan
skenario ini, median jumlah kata turun dari 94 menjadi 60 dan frekuensi LaTeX
\textit{display} memiliki $r=-1{,}00$; satu contoh tersebut dipakai sebagai
ilustrasi bentuk respons, bukan sebagai dasar statistik.

Jumlah kata pada DeepSeek R1-Distill 32B memberi konteks untuk hasil token.
Teks respons pada skenario implisit memuat lebih sedikit kata daripada E,
sedangkan \textit{completion token} meningkat 69,1--74,8\%. Jumlah kata dihitung
dari teks respons, sedangkan \textit{completion token} berasal dari metadata
API. Kedua ukuran tersebut digunakan secara terpisah, sehingga kenaikan token
tidak diartikan sebagai jawaban yang terlihat lebih panjang.

\par\noindent\textbf{d. Pola Pembuka dan Istilah Regional-Kultural}\par

Pola pembuka respons dibandingkan pada N, E, I1, I2, dan I3 untuk setiap model
dan persona. Dalam data regional-kultural, setiap respons masuk ke satu dari
sepuluh kategori yang muncul: deliberatif,
penolakan, jawaban langsung, pembuka berpoin, subjudul, langkah pertama,
aljabar, instruksional, pengulangan soal, atau naratif. Pembuka instruksional
langsung mengajak pengguna menyelesaikan soal, sedangkan naratif langsung masuk
ke isi tanpa pembingkaian yang dikenali aturan kategori. Subjudul dimulai dengan
teks seperti \textbf{Solution}, sedangkan pengulangan soal menuliskan kembali
informasi pada soal.

Contoh pembuka instruksional adalah ``\textit{Let's break it down step by
step}'', contoh subjudul adalah ``\textit{Solution:}'', dan contoh naratif
adalah ``\textit{Geb is 3 years old.}'' Respons Llama 3.1 8B pada
\texttt{gsm8k\_00002} yang dimulai dengan ``\textit{Since the robe takes 2
bolts of blue fiber}'' menjadi contoh pengulangan informasi soal. Kategori
ditetapkan dari bagian awal teks yang terlihat dengan urutan aturan yang sama
pada seluruh model.

Gambar~\ref{fig:regional-opening-gsm8k} memperlihatkan komposisi kesepuluh
kategori. Setiap batang adalah rata-rata setara keenam model untuk satu persona
dan skenario, lalu dinormalisasi menjadi 100\%. Warna menunjukkan kategori
pembuka, bukan tingkat akurasi. Empat kategori dengan proporsi terbesar dibahas
lebih rinci setelah gambar.

\begin{figure}[H]
    \centering
    \includegraphics[
        width=\textwidth,
        height=0.78\textheight,
        keepaspectratio
    ]{image/BabVI/gsm8k/regional_kultural/11_opening_composition_100pct.png}
    \caption{Persentase kategori pembuka respons pada domain regional-kultural}
    \label{fig:regional-opening-gsm8k}
\end{figure}

Pada DeepSeek R1-Distill 32B, sebanyak 61,03\% respons pada N menggunakan
pembuka instruksional dan 38,74\% menggunakan pembuka berupa subjudul. Pada E,
pembuka berupa subjudul mencakup 70,66--79,23\% respons. Pada I1, I2, dan I3,
pembuka naratif mencakup 69,63--73,09\% respons, sedangkan pengulangan soal
mencakup 22,94--25,02\%.

GPT-OSS 20B lebih sering menggunakan pembuka naratif. Persentasenya mencapai
80,91\% pada N, 71,54--74,49\% pada E, dan 61,55--69,89\% pada ketiga skenario
implisit. Pada Qwen3 Next 80B-A3B, pembuka instruksional mencakup 64,72\%
respons pada N dan 47,72--59,97\% pada E, I1, I2, dan I3.

Pada Llama 3.1 8B, Llama 3.3 70B, dan Ministral 8B, pembuka instruksional
mencakup lebih dari 94\% respons pada N. Komposisi E pada Llama 3.1 8B dan
Llama 3.3 70B berubah menurut kawasan, sedangkan Ministral 8B tetap memakai
pembuka instruksional pada 97,57--99,32\% respons E. Rentang pembuka E berbeda
antarmodel dan antarkawasan. Tidak ada satu kategori pembuka yang berubah
dengan pola sama pada seluruh model.

Komposisi pembuka antarkawasan diuji dengan Stuart--Maxwell dan koreksi Holm.
Kategori pembuka diperlakukan sebagai satu hasil multinomial, bukan sepuluh uji
biner yang terpisah.
Hasil signifikan terdapat pada 18 dari 36 pasangan E, sembilan pasangan I1,
serta masing-masing 11 pasangan I2 dan I3.

Pada ketiga skenario implisit, GPT-OSS 20B menghasilkan 11 perbandingan
signifikan. Llama 3.3 70B, Qwen3 Next 80B-A3B, dan Ministral 8B masing-masing
menghasilkan enam, DeepSeek R1-Distill 32B menghasilkan dua, dan Llama 3.1 8B
menghasilkan nol. Hasil setiap pasangan terdapat pada
Lampiran~\ref{fig:lampiran-pembuka-langsung-regional}.

Penyebutan identitas regional-kultural diukur melalui label ``North America'',
``East Asia'', ``South Asia'', dan ``Sub-Saharan Africa''. Label tersebut
dibedakan dari kata seperti ``American'', ``Chinese'', ``Indian'', atau
``African'' yang dapat menjadi bagian dari isi soal GSM8K. Rincian persentase
respons yang memakai label sesuai persona ditempatkan pada
Lampiran~\ref{fig:lampiran-penyebutan-kawasan-regional}. Setiap panel mewakili
satu persona; titik berwarna menunjukkan persentase respons per model dan
skenario, sedangkan berlian hitam menunjukkan median keenam model.

Penyebutan label kawasan hampir seluruhnya muncul pada E dan berbeda
antarmodel. Llama 3.1 8B paling banyak menyebutkan label Asia Timur, Asia
Selatan, dan Afrika Sub-Sahara, sedangkan Llama 3.3 70B menggunakannya dengan
persentase yang lebih rendah. Model lain hampir tidak menyebutkan label kawasan
secara langsung. Label Amerika Utara juga jarang muncul, sehingga penyesuaian
bahasa regional-kultural tidak terjadi secara merata pada empat kawasan.

Istilah regional-kultural yang tercatat juga berbeda menurut kawasan. Pada
Amerika Utara, istilah yang muncul antara lain ``North America'', ``Tradition'',
dan ``Thanksgiving''. Pada Asia Timur, istilah yang muncul mencakup ``East
Asia'', ``Red envelope'', dan ``Hongbao'', sedangkan pada Asia Selatan mencakup
``South Asia'', ``Namaste'', ``Diwali'', dan ``Mehendi''. Pada Afrika
Sub-Sahara, istilah yang muncul mencakup ``Sub-Saharan Africa'', ``Across the
globe'', dan ``Harvest celebration''.

\begin{table}[H]
    \centering
    \caption{Contoh istilah regional-kultural pada skenario persona}
    \label{tab:istilah-regional-gsm8k}
    \footnotesize
    \renewcommand{\arraystretch}{1.15}
    \begin{adjustbox}{max width=\textwidth}
    \begin{tabularx}{\textwidth}{
        >{\raggedright\arraybackslash}p{2.3cm}
        >{\raggedright\arraybackslash}X
        >{\raggedright\arraybackslash}p{1.7cm}
        >{\raggedright\arraybackslash}X
        >{\raggedright\arraybackslash}p{1.9cm}}
        \toprule
        \textbf{Kawasan} & \textbf{E} & \textbf{I1} & \textbf{I2} & \textbf{I3} \\
        \midrule
        Amerika Utara & North America, Tradition & -- & Tradition, Thanksgiving & -- \\
        Asia Timur & East Asia, Culture/cultural & Wei Chen & Red envelope, Hongbao, Culture/cultural & Tradition \\
        Asia Selatan & South Asia, Namaste & -- & Diya, Diwali & Mehendi \\
        Afrika Sub-Sahara & Sub-Saharan Africa, Across the globe & Tradition & Harvest celebration & -- \\
        \bottomrule
    \end{tabularx}
    \end{adjustbox}
\end{table}

Tabel~\ref{tab:istilah-regional-gsm8k} membedakan istilah yang muncul pada E,
I1, I2, dan I3. Label kawasan lebih banyak muncul pada E, sedangkan istilah
pada skenario implisit mengikuti isi \textit{warmup}. I1 dapat membawa nama
``Wei Chen'', I2 dapat membawa perayaan seperti ``Thanksgiving'', ``Hongbao'',
``Diwali'', dan ``Harvest celebration'', sedangkan I3 dapat membawa istilah
``Mehendi''. Soal GSM8K yang digunakan tetap sama, sehingga perbedaan istilah
lebih berkaitan dengan petunjuk persona dan isi \textit{warmup} daripada isi
soal matematika. Daftar istilah hasil audit terdapat pada
Lampiran~\ref{tab:lampiran-istilah-regional}.

Tanda ``--'' pada tabel menandai istilah yang tidak ditemukan dalam ringkasan
audit, bukan bahwa istilah tersebut mustahil muncul. Teks \textit{warmup} tidak
dihitung sebagai respons; istilah baru dihitung jika muncul kembali pada
jawaban utama. Grafik jumlah kemunculan per 10.000 kata ditempatkan pada
Lampiran~\ref{fig:lampiran-istilah-regional-terurut}. Frekuensi dihitung per
model kemudian dirata-ratakan dari keenam model, sehingga frekuensi tinggi pada
satu model tidak dianggap sebagai kebiasaan seluruh model.

Istilah dengan frekuensi tertinggi adalah ``East Asia'' pada persona
Asia Timur, ``South Asia'' dan ``Namaste'' pada persona Asia Selatan, serta
``Sub-Saharan Africa'' pada persona Afrika Sub-Sahara. Istilah Amerika Utara
lebih jarang muncul daripada istilah pada ketiga kawasan tersebut. Prevalensi
sapaan, penyebutan tradisi, dan penyesuaian budaya lainnya terdapat pada
Lampiran~\ref{fig:lampiran-penyesuaian-regional}.

Empat persona di bagian ini adalah proksi rancangan: Anglo Amerika Utara,
Tionghoa Asia Timur, India Asia Selatan, dan Afrika Barat untuk Afrika
Sub-Sahara. Istilah yang terdeteksi karena itu tidak mewakili seluruh budaya di
dalam kawasan tersebut. Penyebutan label, sapaan, atau tradisi juga hanya
menandai penyesuaian yang terlihat dan tidak cukup untuk menyimpulkan kedalaman
pemahaman budaya model.

Audit kontekstual berbantuan AI memeriksa 480 kemunculan penanda linguistik dan tidak
menemukan positif palsu pada sampel tersebut. Pemeriksaan \textit{blind recall}
terhadap 204 respons menemukan tiga respons positif tanpa penanda linguistik yang
terlewat. Karena audit ini belum dilakukan oleh penilai manusia yang independen,
hasil pemeriksaan presisi dan \textit{recall} dibatasi pada sampel audit.

\par\noindent\textbf{e. Penanda Linguistik dan Akurasi}\par

Perubahan bahasa yang paling konsisten terhadap N terdapat pada susunan
respons. Keenam model memakai kalimat yang lebih panjang serta lebih sedikit
langkah eksplisit, daftar berpoin, dan subjudul pada keempat persona dan empat
skenario persona. E juga memunculkan label kawasan dan sapaan regional, tetapi
hasil antarkawasan untuk penanda linguistik khusus hanya terdapat pada Llama 3.1 8B dan
Llama 3.3 70B. I1, I2, dan I3 lebih sering berbeda pada bentuk umum respons;
penanda linguistik khusus regional-kultural pada ketiga skenario ini hanya menghasilkan
tujuh perbandingan signifikan dan semuanya terdapat pada Llama 3.3 70B.

Q Cochran digunakan sebagai uji kelompok. Pada kelompok yang signifikan, Exact
McNemar setelah koreksi Holm menghasilkan sembilan penurunan akurasi I--E.
Tujuh penurunan terdapat pada Qwen3 Next 80B-A3B, sedangkan
DeepSeek R1-Distill 32B dan Ministral 8B masing-masing menghasilkan satu.
Besarnya berkisar 1,06--3,18 poin dengan $|h|=0{,}05$--$0{,}11$. Llama 3.1 8B,
Llama 3.3 70B, dan GPT-OSS 20B tidak menghasilkan perbandingan I--E yang
didukung kedua uji tersebut.

Pada tujuh penurunan Qwen3 Next 80B-A3B, frekuensi kata penghubung sebab-akibat
lebih rendah pada skenario implisit dengan $r=-0{,}36$ sampai $-0{,}25$. Lima
perbandingan Qwen3 Next 80B-A3B juga memiliki frekuensi LaTeX \textit{display}
yang lebih rendah dengan $r=-0{,}27$ sampai $-0{,}15$. Pada DeepSeek
R1-Distill 32B, akurasi I3 persona Amerika Utara lebih rendah 2,05 poin,
jumlah kata turun dari median 94 menjadi 60, dan LaTeX \textit{display}
memiliki $r=-1{,}00$. Pada Ministral 8B, akurasi I2 persona Afrika Sub-Sahara
lebih rendah 3,18 poin dan LaTeX \textit{display} memiliki $r=-0{,}38$, tetapi
jumlah kata naik dari median 106 menjadi 118.

Tujuh penurunan frekuensi LaTeX \textit{display} pada perbandingan akurasi
tetap signifikan ketika analisis dibatasi pada soal yang dijawab benar oleh E
dan skenario implisit. Hasil pemeriksaan kepekaan terdapat pada
Gambar~\ref{fig:lampiran-sensitivitas-penanda-regional}. Perubahan format masih
muncul pada pasangan respons yang sama-sama benar, sehingga penurunan LaTeX
\textit{display} tidak dapat dipakai sebagai penyebab langsung kesalahan.

Perbandingan antarkawasan menghasilkan tiga perbedaan akurasi pada
Exact McNemar setelah koreksi Holm. Pada Ministral 8B I2, akurasi Amerika Utara
dan Asia Selatan lebih tinggi daripada Afrika Sub-Sahara sebesar 4,17 dan 3,41
poin. Pada DeepSeek R1-Distill 32B I3, akurasi Asia Timur lebih tinggi 2,20 poin
daripada Amerika Utara. Perbedaan label kawasan dan sapaan pada E justru
ditemukan pada Llama 3.1 8B dan Llama 3.3 70B. Tidak ada satu kawasan yang
secara konsisten memperoleh akurasi paling tinggi atau paling rendah pada
keenam model.

Analisis memakai pasangan soal dengan teks respons yang tersedia pada kedua
skenario. Terdapat tujuh celah pada berkas sumber N, sedangkan sebagian respons
lain tidak memiliki teks yang dapat dianalisis atau hanya berisi pesan teknis.
Pada 16 skenario persona, ketersediaan teks GPT-OSS 20B paling bervariasi,
yaitu 79,08--97,50\% dari 1.319 soal per skenario; model lain umumnya lebih
lengkap.
Perbedaan ketersediaan ini dicatat pada
Gambar~\ref{fig:lampiran-ketersediaan-penanda-regional} dan ukuran pasangan
efektif ikut dilaporkan pada keluaran statistik.

Persona regional-kultural lebih konsisten mengubah cara jawaban ditulis
daripada akurasinya. Perubahan yang konsisten lintas model
terdapat pada panjang kalimat dan pengurangan penanda linguistik struktur, bukan pada satu
kawasan tertentu. Label dan sapaan lebih bergantung pada model, sedangkan
perbedaan akurasi hanya muncul pada model, persona, dan skenario tertentu.
Rancangan ini tidak memakai \textit{warmup} netral, sehingga perbandingan I--N
masih mencakup isi regional-kultural sekaligus penambahan riwayat percakapan.

\subsubsection{Ringkasan Domain Regional-Kultural}
\label{subsubsec:ringkasan-regional-gsm8k}

Pada domain regional-kultural, bagian respons yang paling banyak berubah
berbeda menurut cara persona disampaikan. E lebih banyak memunculkan label
kawasan dan sapaan regional, terutama pada Llama 3.1 8B dan Llama 3.3 70B. Pada
I1, I2, dan I3, istilah regional-kultural lebih sedikit dan perubahan lebih
banyak terdapat pada susunan jawaban setelah riwayat percakapan ditambahkan.

Untuk rumusan masalah pertama, persona regional-kultural lebih konsisten
mengubah penanda linguistik daripada akurasi respons. Dibandingkan dengan N,
perubahan penanda linguistik ditemukan pada keenam model. Perbandingan I--E
yang didukung Q Cochran dan koreksi Holm menghasilkan sembilan penurunan
akurasi, yaitu tujuh pada Qwen3 Next 80B-A3B serta masing-masing satu pada
DeepSeek R1-Distill 32B dan Ministral 8B. Penurunan tersebut muncul pada persona
Amerika Utara, Asia Selatan, dan Afrika Sub-Sahara, sedangkan persona Asia Timur
tidak menghasilkan penurunan I--E yang didukung kedua uji. Perbandingan
antarkawasan hanya menghasilkan tiga perbedaan akurasi dan tidak menunjukkan
satu kawasan yang selalu memperoleh akurasi paling tinggi atau paling rendah.

Untuk bagian rumusan masalah kedua yang berkaitan dengan tugas numerik,
skenario implisit menggunakan total token 79,1--330,1\% lebih banyak daripada E
karena riwayat percakapan ikut dikirim bersama soal. Perubahan \textit{completion
token} berbeda antarmodel dan tidak selalu diikuti perubahan akurasi yang sama.
Persona regional-kultural dapat mengubah bentuk respons, penggunaan
token, dan akurasi dengan pola yang berbeda pada tugas GSM8K.

Untuk rumusan masalah ketiga, pengaruh persona regional-kultural tidak
konsisten antarmodel. Perbedaan label kawasan dan sapaan paling banyak muncul
pada dua model Llama, sedangkan penurunan akurasi muncul pada Qwen3 Next
80B-A3B, DeepSeek R1-Distill 32B, dan Ministral 8B. Llama 3.1 8B dan Ministral
8B memiliki ukuran parameter yang sama, tetapi menghasilkan bentuk pengaruh
yang berbeda. Ukuran parameter saja tidak menghasilkan pola pengaruh persona
yang sama.

% =====================================================================
\subsection{Perbandingan Antardomain Persona}
\label{subsec:perbandingan-domain-gsm8k}

Hasil GSM8K pada domain gender, agama, dan regional-kultural dibandingkan untuk
melihat perbedaan pengaruh persona pada tugas numerik. Gambar~\ref{fig:omnibus-semua-domain-gsm8k}
menunjukkan hasil Q Cochran untuk setiap model dan persona. Uji tersebut
menentukan kelompok yang memiliki perbedaan akurasi antarskenario, tetapi tidak
menentukan pasangan skenario yang berbeda. Pasangan skenario dan tanda
selisihnya ditentukan melalui Exact McNemar setelah koreksi Holm.

\begin{landscape}
\begin{figure}[H]
    \centering
    \includegraphics[
        width=\linewidth,
        height=0.72\textheight,
        keepaspectratio
    ]{image/bab6/gsm8k/semua_domain/akurasi/20_all_groups_cochran_q_omnibus_pvalue_heatmap.png}
    \caption{Nilai-p uji omnibus Q Cochran pada seluruh domain persona GSM8K}
    \label{fig:omnibus-semua-domain-gsm8k}
\end{figure}
\end{landscape}

Tanda bintang pada Gambar~\ref{fig:omnibus-semua-domain-gsm8k} menandai
$p<0{,}05$. Pada gender, Q Cochran membandingkan N, E, I1, I2, dan I3, sedangkan
pada agama dan regional-kultural uji tersebut membandingkan E, I1, I2, dan I3.
Perbedaan ini terjadi karena N pada gender menggunakan protokol inferensi yang
sama, sedangkan N pada agama dan regional-kultural menjadi analisis sensitivitas
karena pelaksanaannya tidak sepenuhnya sama. Nilai $p$ pada gambar tidak
disetarakan karena jumlah skenario dalam uji berbeda.
Tabel~\ref{tab:identitas-terdampak-gsm8k} dan
Tabel~\ref{tab:model-jenis-ukuran-gsm8k} hanya menggunakan pasangan I--E agar
jenis perbandingannya sama.

Pengaruh persona pada akurasi berbeda menurut domain dan model. Q Cochran
signifikan pada tiga dari 12 kelompok gender, delapan dari 24 kelompok agama,
dan enam dari 24 kelompok regional-kultural. Dua kelompok gender terdapat pada
DeepSeek R1-Distill 32B dan satu kelompok terdapat pada Llama 3.3 70B, sedangkan
empat kelompok agama berasal dari seluruh persona pada Ministral 8B. Pada
regional-kultural, tiga kelompok terdapat pada Qwen3 Next 80B-A3B dan tiga
lainnya masing-masing terdapat pada Llama 3.1 8B, DeepSeek R1-Distill 32B, dan
Ministral 8B.

Perbedaan penanda linguistik ditemukan lebih luas daripada perbedaan akurasi.
Pada agama dan regional-kultural, E lebih sering memunculkan label identitas
dan istilah yang sesuai dengan persona. Pola tersebut sesuai dengan E yang
menyebut identitas secara langsung, sedangkan I1--I3 menambahkan percakapan
\textit{warmup} tanpa mengubah soal matematika. Pada gender, perubahan lebih
banyak terdapat pada pronomina, bahasa matematis, dan pola pembuka respons.
Perubahan penanda linguistik dan akurasi tidak selalu muncul pada model yang
sama sehingga kemunculan istilah persona tidak langsung menunjukkan penurunan
akurasi.

Penggunaan total token meningkat pada skenario implisit di ketiga domain.
Rentang kenaikannya terhadap E adalah 59,5--220,5\% pada gender,
129,2--409,5\% pada agama, dan 79,1--330,1\% pada regional-kultural. Kenaikan
tersebut terutama berasal dari \textit{input token} karena percakapan
\textit{warmup} ikut dikirim bersama setiap soal. Perubahan \textit{completion
token} berbeda antarmodel dan tidak selalu diikuti perubahan akurasi yang sama. Pada
layanan yang menghitung biaya berdasarkan jumlah token, kenaikan \textit{input
token} menambah komponen biaya input, sedangkan total biayanya mengikuti harga
\textit{input} dan \textit{output token}\label{term:output-token} pada setiap model.

Tabel~\ref{tab:perbandingan-domain-gsm8k} merangkum hasil Holm menurut jenis
perbandingan. Pada agama dan regional-kultural, penyebut 96 untuk N berasal dari
6 model $\times$ 4 persona $\times$ 4 skenario, sedangkan penyebut 72 untuk
I--E berasal dari 6 model $\times$ 4 persona $\times$ 3 skenario implisit.
Penyebut 144 pada perbandingan antarpersona berasal dari 6 model $\times$ 4
skenario $\times$ 6 pasangan persona. Pada gender, perhitungan yang sama
menggunakan dua persona sehingga menghasilkan 48, 36, dan 24 perbandingan.

\begin{center}
    \centering
    \captionof{table}{Ringkasan perbandingan hasil GSM8K antardomain persona}
    \label{tab:perbandingan-domain-gsm8k}
    \footnotesize
    \renewcommand{\arraystretch}{1.16}
    \begin{tabularx}{\textwidth}{
        >{\raggedright\arraybackslash}p{1.8cm}
        >{\centering\arraybackslash}p{2.8cm}
        >{\raggedright\arraybackslash}X
        >{\raggedright\arraybackslash}X}
        \toprule
        \textbf{Domain} &
        \textbf{Hasil setelah Holm} &
        \textbf{Pola akurasi} &
        \textbf{Pola penanda linguistik} \\
        \midrule
        Gender &
        N: 4/48; I--E: 4/36; antarpersona: 0/24 &
        Tiga penurunan I--E terdapat pada DeepSeek R1-Distill 32B dan satu
        kenaikan terdapat pada Llama 3.3 70B. &
        Perubahan terlihat pada pronomina gender, bahasa matematis, dan pola
        pembuka respons. \\
        \addlinespace[3pt]

        Agama &
        N: 23/96; I--E: 19/72; antarpersona: 4/144 &
        Perubahan I--E terbesar terdapat pada persona Islam di Ministral 8B.
        Keempat perbandingan antarpersona melibatkan persona Islam dengan akurasi
        yang lebih rendah. &
        E lebih banyak memunculkan identitas dan istilah agama, sedangkan
        I1--I3 lebih banyak mengubah pola respons model. \\
        \addlinespace[3pt]

        Regional-kultural &
        N: 15/96; I--E: 10/72; antarpersona: 3/144 &
        Sembilan hasil I--E didukung Q Cochran dan Holm; seluruhnya berupa
        penurunan pada Amerika Utara, Asia Selatan, atau Afrika Sub-Sahara. &
        Perubahan terlihat pada penyebutan kawasan, istilah budaya, sapaan,
        tradisi atau perayaan, dan kategori pembuka respons. \\
        \bottomrule
    \end{tabularx}
\end{center}

Hasil pada Tabel~\ref{tab:perbandingan-domain-gsm8k} menunjukkan bahwa persona
dapat mengubah akurasi dan pola respons model, tetapi bentuknya berbeda pada
setiap domain. Agama memiliki perubahan I--E terbesar dan seluruh perbandingan
antarpersona yang signifikan melibatkan persona Islam dengan akurasi lebih rendah.
Gender tidak memiliki perbedaan antarpersona antara persona laki-laki dan
perempuan, sedangkan tiga perbedaan antarkawasan regional-kultural hanya muncul
pada I2 atau I3 di Ministral 8B dan DeepSeek R1-Distill 32B.

\subsubsection{Perbandingan Akurasi Skenario Implisit dan Eksplisit per Persona}
\label{subsubsec:identitas-terdampak-gsm8k}

Tabel~\ref{tab:identitas-terdampak-gsm8k} memuat perbandingan I--E yang didukung
Q Cochran dan koreksi Holm. Hasil ini menunjukkan persona yang akurasinya
berubah ketika cara penyampaian persona diubah dari eksplisit menjadi implisit.
Tabel tersebut tidak membandingkan akurasi satu persona dengan persona lain.
Persona yang tidak tercantum tidak memiliki hasil I--E yang didukung kedua uji.

\begin{table}[htbp]
    \centering
    \caption{Persona dengan perbedaan akurasi I--E yang didukung Q Cochran dan koreksi Holm}
    \label{tab:identitas-terdampak-gsm8k}
    \footnotesize
    \renewcommand{\arraystretch}{1.16}
    \begin{tabularx}{\textwidth}{
        >{\raggedright\arraybackslash}p{1.9cm}
        >{\raggedright\arraybackslash}p{2.6cm}
        >{\centering\arraybackslash}p{1.6cm}
        >{\centering\arraybackslash}p{2.3cm}
        >{\raggedright\arraybackslash}X}
        \toprule
        \textbf{Domain} & \textbf{Persona} & \textbf{Perubahan} &
        \textbf{Rentang $h$} & \textbf{Model} \\
        \midrule
        Gender & Laki-laki dan perempuan & Turun & $-0{,}12$ s.d. $-0{,}09$ &
        DeepSeek R1-Distill 32B \\
        Gender & Perempuan (I3) & Naik & $+0{,}08$ & Llama 3.3 70B \\
        \addlinespace[3pt]
        Agama & Islam & Turun paling besar & $-0{,}27$ s.d. $-0{,}15$ &
        Ministral 8B \\
        Agama & Islam & Naik & $+0{,}07$ s.d. $+0{,}10$ &
        GPT-OSS 20B \\
        Agama & Ateis dan Yahudi & Turun & $-0{,}13$ s.d. $-0{,}09$ &
        Ministral 8B \\
        Agama & Kristen & Turun & $-0{,}16$ s.d. $-0{,}06$ &
        Ministral 8B, Llama 3.1 8B, Qwen3 Next 80B-A3B \\
        \addlinespace[3pt]
        Regional & Asia Selatan & Turun & $-0{,}07$ s.d. $-0{,}06$ &
        Qwen3 Next 80B-A3B \\
        Regional & Afrika Sub-Sahara & Turun & $-0{,}11$ s.d. $-0{,}05$ &
        Qwen3 Next 80B-A3B, Ministral 8B \\
        Regional & Amerika Utara & Turun & $-0{,}08$ s.d. $-0{,}07$ &
        Qwen3 Next 80B-A3B, DeepSeek R1-Distill 32B \\
        Regional & Asia Timur & Tidak ada & --- & --- \\
        \bottomrule
    \end{tabularx}
\end{table}

Perubahan I--E terbesar terdapat pada persona Islam. Pada Ministral 8B,
akurasi I1 persona Islam turun 8,79 poin dari E dengan $h=-0{,}27$, sedangkan
GPT-OSS 20B menghasilkan kenaikan 2,06--2,67 poin pada I1--I3. Pada gender,
penurunan terdapat pada persona laki-laki dan perempuan di DeepSeek R1-Distill
32B, sedangkan I3 persona perempuan meningkat pada Llama 3.3 70B. Pada
regional-kultural, Amerika Utara, Asia Selatan, dan Afrika Sub-Sahara
masing-masing memiliki tiga penurunan yang didukung kedua uji. Asia Timur tidak
memiliki hasil I--E yang didukung kedua uji.

Perbandingan antarpersona memberi batas yang berbeda dari hasil I--E.
Keempat hasil agama melibatkan persona Islam dengan akurasi lebih rendah pada
Ministral 8B atau Llama 3.1 8B. Pada regional-kultural, Afrika Sub-Sahara lebih
rendah dalam dua perbandingan pada Ministral 8B, sedangkan Amerika Utara lebih
rendah daripada Asia Timur dalam satu perbandingan pada DeepSeek R1-Distill
32B. Gender tidak memiliki perbedaan antarpersona yang signifikan sehingga hasilnya
tidak mendukung penetapan persona laki-laki atau perempuan sebagai persona yang
lebih rentan pada GSM8K.

\subsubsection{Perbandingan Berdasarkan Jenis dan Ukuran Model}
\label{subsubsec:jenis-ukuran-model-gsm8k}

Pertanyaan penelitian ketiga membahas pengaruh persona pada model dengan
jenis dan ukuran yang berbeda. Tabel~\ref{tab:model-jenis-ukuran-gsm8k}
merangkum perbandingan I--E yang didukung Q Cochran dan koreksi Holm pada setiap
model. Perbandingan terhadap skenario netral dan perbandingan antarpersona tidak
disertakan agar ketiga domain menggunakan jenis perbandingan yang sama.
Pengelompokan ukuran mengikuti jumlah parameter yang dilaporkan pada
Bab~\ref{chap:perancangan}.

\begin{center}
    \centering
    \captionof{table}{Jumlah perbandingan I--E yang didukung Q Cochran dan koreksi Holm per model}
    \label{tab:model-jenis-ukuran-gsm8k}
    \small
    \renewcommand{\arraystretch}{1.16}
    \begin{tabularx}{\textwidth}{
        >{\raggedright\arraybackslash}X
        >{\centering\arraybackslash}p{2.1cm}
        >{\centering\arraybackslash}p{2.3cm}
        >{\centering\arraybackslash}p{1.3cm}
        >{\centering\arraybackslash}p{1.3cm}
        >{\centering\arraybackslash}p{1.5cm}
        >{\centering\arraybackslash}p{1.1cm}}
        \toprule
        \textbf{Model} & \textbf{Jenis} & \textbf{Ukuran} &
        \textbf{Gender} & \textbf{Agama} & \textbf{Regional} &
        \textbf{Total} \\
        \midrule
        Llama 3.1 8B & Non-reasoning & Kecil (8B) & 0 & 2 & 0 & 2 \\
        Ministral 8B & Non-reasoning & Kecil (8B) & 0 & 12 & 1 & 13 \\
        DeepSeek R1-Distill 32B & Reasoning & Menengah (32B) & 3 & 0 & 1 & 4 \\
        GPT-OSS 20B & Reasoning & Menengah (20B) & 0 & 3 & 0 & 3 \\
        Llama 3.3 70B & Non-reasoning & Besar (70B) & 1 & 0 & 0 & 1 \\
        Qwen3 Next 80B-A3B & Reasoning & Besar (80B, MoE) & 0 & 2 & 7 & 9 \\
        \midrule
        \multicolumn{3}{r}{\textbf{Total}} & 4 & 19 & 9 & 32 \\
        \bottomrule
    \end{tabularx}
\end{center}

Model \textit{reasoning} dan \textit{non-reasoning} masing-masing menghasilkan
16 dari 32 hasil I--E. Model kecil menghasilkan 15 hasil, model menengah tujuh
hasil, dan model besar 10 hasil sehingga jumlahnya tidak membentuk urutan
berdasarkan ukuran. Llama 3.1 8B dan Ministral 8B juga menghasilkan dua dan 13
hasil walaupun ukuran parameternya sama. Jenis dan ukuran parameter karena itu
tidak membentuk pola pengaruh persona yang sama pada keenam model.

Qwen3 Next 80B-A3B dan GPT-OSS 20B menggunakan arsitektur
\textit{mixture-of-experts}. Jumlah parameter total pada arsitektur tersebut
berbeda dari jumlah parameter yang aktif pada setiap token. Pengelompokan ukuran
di tabel karena itu tidak mencakup seluruh perbedaan komputasi dan arsitektur
antarmodel.

% =====================================================================
\subsection{Ringkasan Evaluasi GSM8K}
\label{subsec:ringkasan-evaluasi-gsm8k}

Evaluasi GSM8K menunjukkan bahwa pengaruh persona berbeda menurut domain,
model, dan cara persona disampaikan. Pada gender, empat perbandingan I--E
didukung Q Cochran dan Holm, tetapi tidak ada perbedaan antarpersona antara persona
laki-laki dan perempuan. Pada agama, 19 perbandingan I--E didukung kedua uji dan
penurunan terbesar terdapat pada persona Islam di Ministral 8B. Keempat
perbandingan antarpersona agama juga melibatkan persona Islam dengan akurasi lebih
rendah. Pada regional-kultural, sembilan perbandingan I--E yang didukung kedua
uji seluruhnya berupa penurunan pada Amerika Utara, Asia Selatan, atau Afrika
Sub-Sahara; dua dari tiga perbandingan antarkawasan menempatkan Afrika Sub-Sahara
lebih rendah.

Perubahan penanda linguistik ditemukan lebih luas daripada perubahan akurasi.
Skenario eksplisit lebih sering memunculkan label identitas atau istilah agama
dan budaya, sedangkan skenario implisit lebih banyak mengubah susunan respons.
Pada gender, perubahan juga terdapat pada pronomina, bahasa matematis, dan pola
pembuka. Total token skenario implisit meningkat 59,5--220,5\% pada gender,
129,2--409,5\% pada agama, dan 79,1--330,1\% pada regional-kultural. Kenaikan
konteks input menambah komponen biaya input, tetapi perubahan token tidak selalu
diikuti perubahan akurasi yang sama.

Jenis dan ukuran parameter tidak menghasilkan pola pengaruh persona yang sama.
Model \textit{reasoning} dan \textit{non-reasoning} masing-masing menghasilkan
16 dari 32 perbandingan I--E yang didukung kedua uji. Model kecil menghasilkan
15 hasil, model menengah tujuh hasil, dan model besar 10 hasil. Llama 3.1 8B dan
Ministral 8B memiliki ukuran yang sama, tetapi masing-masing menghasilkan dua
dan 13 hasil.

% =====================================================================
\section{Evaluasi ETHICS}
\label{sec:evaluasi-ethics}

ETHICS terdiri atas 1.514 soal dari subset \textit{Commonsense},
\textit{Deontology}, \textit{Justice}, dan \textit{Virtue}. Pengujian mencakup
satu skenario netral dan 40 skenario persona untuk setiap model, sesuai matriks
pengujian pada Subbab~\ref{subsec:matriks-pengujian}. Perbandingan
antarskenario mengikuti rancangan analisis statistik pada
Subbab~\ref{sec:rancangan-analisis-statistik}. Skenario netral merupakan
skenario tanpa persona yang digunakan sebagai \textit{baseline}.

Pada pembaruan domain gender dan regional-kultural, akurasi dihitung pada respons
\textit{committed} untuk \textit{Commonsense} dan \textit{Justice}.
\textit{Deontology} dan \textit{Virtue} dianalisis melalui distribusi
keputusan model karena pertanyaan yang dikirim tidak sesuai dengan tugas
penentu label dataset. Proporsi \textit{committed}, \textit{abstain}, dan
\textit{missing} dicantumkan terpisah, sedangkan keberhasilan efektif
memasukkan respons \textit{abstain} ke dalam penyebut respons nonmissing.
Respons \textit{missing} tidak ditafsirkan sebagai perilaku moral model karena
teks respons tidak tersedia.

Hasil agama masih menggunakan analisis sebelumnya,
dengan seluruh respons sebagai penyebut kecocokan terhadap label dataset.
Angka pada \textit{Deontology} dan \textit{Virtue} dalam hasil agama tersebut
tidak digunakan untuk menyimpulkan ketepatan moral terhadap tugas asli.
Perbandingan proporsi hasil signifikan yang melibatkan domain agama ditunda
sampai subset, penyebut, dan aturan pemilihan pasangan disamakan.

% =====================================================================
\subsection{Domain Gender}
\label{subsec:hasil-gender-ethics}

Domain gender mencakup persona laki-laki dan perempuan. Setiap persona diuji
dengan satu skenario eksplisit dan tiga skenario implisit. I1 menggunakan peran
dalam keluarga, I2 menggunakan kegiatan merapikan penampilan, dan I3
menggunakan pilihan pakaian. Ketiganya dianalisis secara terpisah karena
menggunakan \textit{warmup} yang berbeda. Analisis ini menggunakan 81.756
respons dari 1.514 soal ETHICS.

\subsubsection{Akurasi, Keputusan, dan Status Respons}
\label{subsubsec:akurasi-gender-ethics}

\input{Bab VI - ETHICS Gender Akurasi.tex}

\par\noindent\textbf{e. Penggunaan Token}\par
\label{subsubsec:token-gender-ethics}

Analisis token menggunakan 80.919 dari 81.756 respons ETHICS gender karena 837
respons tidak memiliki metadata token. Rata-rata token dihitung menurut model,
subset, persona, dan skenario. E dibandingkan dengan I1--I3 pada 36 pasangan
per subset, sedangkan E/I1/I2/I3 dibandingkan dengan N pada 48 pasangan per
subset. Metadata yang digunakan mencakup \textit{input token},
\textit{completion token}, \textit{reasoning token}, dan total token.

Pada perbandingan I1--I3 dengan E, \textit{input token} meningkat pada seluruh
36 pasangan di setiap subset. Rentangnya adalah $+42{,}74$--$+697{,}02$\% pada
\textit{Commonsense}, $+151{,}48$--$+3.458{,}00$\% pada \textit{Deontology},
$+171{,}30$--$+2.788{,}41$\% pada \textit{Justice}, dan
$+175{,}57$--$+2.820{,}35$\% pada \textit{Virtue}. Percakapan
\textit{warmup} ikut dikirim bersama soal pada I1--I3 sehingga menambah konteks
input dibandingkan dengan E.

\textit{Completion token} pada I--E meningkat/menurun sebanyak 21/15, 20/16,
23/13, dan 17/19 pada empat subset. \textit{Reasoning token} tercatat pada 18
pasangan per subset dengan jumlah kenaikan/penurunan 3/15, 2/16, 3/15, dan
3/15. Total token meningkat/menurun sebanyak 34/2, 26/10, 31/5, dan 31/5.
Rentang total token di seluruh subset adalah $-36{,}18$ sampai $+510{,}46$\%.

Pada perbandingan skenario persona dengan N, \textit{input token} meningkat
pada seluruh 48 pasangan setiap subset. \textit{Completion token}
meningkat/menurun sebanyak 21/27, 14/34, 23/25, dan 18/30, sedangkan total token
meningkat/menurun sebanyak 36/12, 28/20, 38/10, dan 36/12. \textit{Reasoning
token} tercatat pada 24 pasangan per subset dengan jumlah kenaikan/penurunan
6/18, 4/20, 7/17, dan 6/18.

Perubahan token dipisahkan dari perubahan akurasi dan status respons. Kenaikan
\textit{input token} tidak selalu disertai kenaikan \textit{completion token},
\textit{reasoning token}, total token, atau akurasi. Biaya nominal juga tidak
dihitung karena harga token input dan output berbeda menurut model dan
penyedia. Ringkasan per subset terdapat pada
Lampiran~\ref{tab:token-gender-ethics-ringkas}.

Visualisasi rinci rata-rata, perubahan terhadap netral, dan hubungan
token--akurasi ditempatkan sebagai materi pendukung agar bagian utama tetap
berfokus pada keluaran ETHICS.

\subsubsection{Ringkasan Akurasi dan Keputusan Domain Gender}
\label{subsubsec:ringkasan-akurasi-gender-ethics}

Empat dari 48 perbandingan akurasi antara persona laki-laki dan perempuan
signifikan setelah koreksi Holm. Seluruhnya terdapat pada \textit{Justice},
dengan tiga hasil lebih tinggi pada persona laki-laki dan satu pada persona
perempuan. Perbandingan ini hanya menggunakan pasangan soal yang menghasilkan
keputusan pada kedua skenario. \textit{Commonsense} tidak menghasilkan
perbandingan akurasi yang signifikan pada ketiga keluarga pengujian.

Pada perbandingan antarskenario, enam dari 72 hasil I--E dan dua dari 96
hasil persona--N didukung Q Cochran serta koreksi Holm.
Seluruhnya berupa penurunan akurasi \textit{Justice}, dengan enam hasil pada
Qwen3 Next 80B-A3B dan dua pada Ministral 8B.
Jumlah pasangan yang masuk ke setiap uji berbeda karena tidak semua skenario
menghasilkan keputusan pada soal yang sama. Hasil tersebut tidak
digeneralisasi ke respons yang berstatus \textit{abstain} atau
\textit{missing}.

Gabungan empat subset menunjukkan perubahan distribusi keputusan, bukan
akurasi moral gabungan. Proporsi \textit{wrong} pada gabungan skenario
implisit lebih tinggi daripada E pada Llama 3.1 8B, tetapi lebih rendah pada
DeepSeek R1-Distill 32B untuk kedua persona. Sebanyak 76,68\% respons
berstatus \textit{committed}, 22,65\% \textit{abstain}, dan 0,67\%
\textit{missing}. Distribusi keputusan dan frekuensi model menghasilkan
keputusan karena itu dipisahkan dari ketepatan terhadap label dataset.

\subsubsection{Penanda Linguistik}
\label{subsubsec:penanda-gender-ethics}

Penanda linguistik ETHICS domain gender dihitung melalui notebook v1.5
\textit{candidate} yang terpisah dari analisis akurasi. Bagian akurasi tidak
menggunakan hasil notebook tersebut untuk menentukan benar--salah, status
respons, atau signifikansi perbandingan akurasi. Hasil penanda linguistik tetap
perlu dilaporkan dengan status \textit{candidate} sampai validasi manusia
independen selesai.

\input{Bab VI - ETHICS Gender Penanda Linguistik.tex}

\subsubsection{Ringkasan Domain Gender}
\label{subsubsec:ringkasan-gender-ethics}

Pada GSM8K, tidak ada satu pun dari 24 perbandingan akurasi antarpersona
laki-laki--perempuan yang signifikan setelah koreksi Holm. Pada
\textit{Commonsense} dan \textit{Justice} ETHICS, empat dari 48 perbandingan
akurasi pada pasangan \textit{committed} signifikan.
Jumlah hasil signifikan tidak digunakan untuk menyimpulkan bahwa pengaruh
persona lebih besar pada salah satu tugas karena ukuran, jumlah soal, dan
pasangan yang dianalisis berbeda.

Pengaruh persona berbeda antarmodel. Qwen3 Next 80B-A3B menghasilkan enam
perbedaan akurasi antarskenario pada \textit{Justice}, tetapi tidak memiliki
perbedaan akurasi laki-laki--perempuan yang signifikan.
GPT-OSS 20B justru memiliki dua perbedaan laki-laki--perempuan.
Jenis maupun ukuran model tidak digunakan sebagai penyebab tunggal
perbedaan respons tersebut.

% =====================================================================
\subsection{Domain Agama}
\label{subsec:hasil-agama-ethics}

Domain agama mencakup persona Islam, Kristen, Yahudi, dan ateis. Setiap persona
diuji dengan skenario E, I1, I2, dan I3, sedangkan N digunakan sebagai
pembanding tanpa persona. Keempat subset memiliki jumlah soal berbeda dan
karena itu tetap dipisahkan pada seluruh pengujian. Analisis akurasi memakai
154.428 respons.

\subsubsection{Akurasi Respons}
\label{subsubsec:akurasi-agama-ethics}

\input{Bab VI - ETHICS Religion Akurasi.tex}

\subsubsection{Ringkasan Akurasi Domain Agama}
\label{subsubsec:ringkasan-akurasi-agama-ethics}

\textit{Virtue} paling sering menghasilkan perbedaan akurasi, dengan 103
hasil antarskenario utama dan 41 hasil antarpersona agama. \textit{Justice}
menghasilkan 93 hasil antarskenario dan 30 hasil antarpersona, sedangkan
\textit{Commonsense} dan \textit{Deontology} masing-masing menghasilkan 83
dan 86 hasil antarskenario. Qwen3 Next 80B-A3B paling sering berubah pada
perbandingan antarskenario, yaitu 95 hasil. DeepSeek R1-Distill 32B paling
banyak menghasilkan perbedaan antarpersona agama, yaitu 24 hasil.

Perubahan I--E terdiri atas 88 kenaikan dan 98 penurunan, sedangkan
perbandingan persona dengan N terdiri atas 40 kenaikan dan 139 penurunan.
Persona Islam paling sering berada pada sisi dengan akurasi lebih tinggi dalam
hasil antarpersona agama, yaitu 35 hasil, dan berada pada sisi lebih rendah
pada 14 hasil. Kecenderungan ini tidak konsisten pada seluruh model, subset,
dan skenario. Hasil tersebut juga tidak menunjukkan kemampuan moral suatu
kelompok agama.

Perubahan status respons menyertai sebagian besar hasil akurasi. Pada 359 dari
365 hasil antarskenario utama, sisi dengan akurasi penuh lebih tinggi memiliki
\textit{abstain} lebih rendah; 261 hasil juga memiliki akurasi
\textit{committed} lebih tinggi. E berbeda dari I1--I3 pada 186 dari 288
perbandingan, dengan 88 kenaikan dan 98 penurunan, sedangkan perbandingan dengan
N lebih sering menurun. Angka analisis sebelumnya ini tidak dibandingkan
dengan hasil gender yang sudah menggunakan akurasi pada pasangan
\textit{committed} untuk dua subset.

Pada GSM8K domain agama, empat dari 144 perbandingan antarpersona agama
signifikan. Seluruhnya melibatkan persona Islam pada Ministral 8B atau
Llama 3.1 8B. Perbedaan jumlah hasil signifikan antartugas tidak digunakan
sebagai ukuran besar pengaruh persona.

% =====================================================================
\subsection{Domain Regional-Kultural}
\label{subsec:hasil-regional-ethics}

Domain regional-kultural mencakup persona Amerika Utara, Asia Timur, Asia
Selatan, dan Afrika Sub-Sahara. Skenario E, I1, I2, dan I3 diperiksa terhadap
N dan dibandingkan pada persona kawasan yang sama. Perbandingan antarkawasan
menggunakan soal, model, subset, dan skenario yang berpasangan. Analisis
akurasi memakai 154.428 respons.

\subsubsection{Akurasi Respons}
\label{subsubsec:akurasi-regional-ethics}

\input{Bab VI - ETHICS Regional Akurasi.tex}

\subsubsection{Ringkasan Akurasi Domain Regional-Kultural}
\label{subsubsec:ringkasan-akurasi-regional-ethics}

Pada perbandingan setiap persona dengan N, perubahan \textit{coverage}
lebih sering muncul daripada perubahan akurasi respons yang sudah
\textit{committed}: 178 dari 384 perbandingan \textit{coverage} signifikan
setelah Holm, sedangkan akurasi \textit{committed} hanya sembilan dari 192
perbandingan pada \textit{Commonsense} dan \textit{Justice}. Keberhasilan
efektif, yang memeriksa ketepatan sekaligus ketersediaan putusan, signifikan
pada 68 dari 192 perbandingan. Pada \textit{Deontology} dan \textit{Virtue},
42 dari 192 perbandingan distribusi keputusan signifikan setelah Holm.

Hasil tersebut berubah menurut model, persona, subset, dan skenario sehingga
tidak membentuk urutan kawasan yang umum. Interpretasi tetap dilakukan per
subset dan keluaran; angka tersebut tidak digunakan untuk menyatakan bahwa
satu kawasan lebih akurat atau lebih bermoral. Rincian kombinasi model,
persona, skenario, dan subset ditampilkan pada lampiran tanpa membentuk satu
angka ``akurasi ETHICS gabungan''.

Pada GSM8K regional-kultural, sembilan dari 72 hasil I--E didukung Q Cochran
dan Holm serta tiga dari 144 perbandingan antarkawasan tetap signifikan.
Perbedaan jumlah hasil signifikan antara GSM8K dan ETHICS tidak digunakan
sebagai ukuran besar pengaruh persona karena tugas, jumlah soal, keluaran yang
dinilai, dan keluarga perbandingannya berbeda.

% =====================================================================
\subsection{Ringkasan Status Respons dan Batas Perbandingan ETHICS Antardomain}
\label{subsec:ringkasan-akurasi-antardomain-ethics}
\label{acr:rq}

Perbandingan proporsi hasil akurasi yang signifikan dengan domain agama belum
dilakukan karena ukuran agama masih memakai seluruh respons pada empat subset.
Gender dan regional-kultural sama-sama menggunakan akurasi pasangan
\textit{committed} pada \textit{Commonsense} dan \textit{Justice}, serta
memisahkan \textit{coverage} dan distribusi keputusan pada dua subset lain.
Jumlah hasil signifikan keduanya tetap tidak dipakai sebagai ukuran besar
pengaruh persona karena jumlah persona dan keluarga perbandingannya berbeda.
Angka kecocokan label pada \textit{Deontology} dan \textit{Virtue} tidak
digunakan untuk menyimpulkan akurasi moral terhadap tugas aslinya.

Tabel~\ref{tab:ringkasan-akurasi-antardomain-ethics} membandingkan status
respons, yang pada ketiga domain menggunakan seluruh respons sebagai
penyebut. \textit{Committed} menunjukkan keputusan yang diperoleh aturan
ekstraksi, sedangkan \textit{abstain} menunjukkan teks tanpa keputusan yang
dapat digunakan. \textit{Missing} dicantumkan terpisah karena menunjukkan
respons yang tidak tersedia, bukan keputusan moral model.

\begin{table}[H]
    \centering
    \caption{Status respons ETHICS menurut domain (\%)}
    \label{tab:ringkasan-akurasi-antardomain-ethics}
    \small
    \renewcommand{\arraystretch}{1.16}
    \begin{tabularx}{\textwidth}{Xccc}
        \toprule
        \textbf{Domain} & \textbf{\textit{Committed}} &
        \textbf{\textit{Abstain}} & \textbf{\textit{Missing}} \\
        \midrule
        Gender & 76,68 & 22,65 & 0,67 \\
        Agama & 68,63 & 29,83 & 1,54 \\
        Regional-kultural & 72,55 & 26,27 & 1,18 \\
        \bottomrule
    \end{tabularx}
\end{table}

Gender memiliki proporsi \textit{committed} tertinggi, yaitu 76,68\%,
dibandingkan 72,55\% pada regional-kultural dan 68,63\% pada agama.
Agama memiliki proporsi \textit{abstain} dan \textit{missing} tertinggi,
masing-masing 29,83\% dan 1,54\%.
Perbedaan tersebut menunjukkan frekuensi model menghasilkan keputusan dan
ketersediaan respons, bukan keunggulan akurasi moral suatu domain.
